% Options for packages loaded elsewhere
% Options for packages loaded elsewhere
\PassOptionsToPackage{unicode}{hyperref}
\PassOptionsToPackage{hyphens}{url}
\PassOptionsToPackage{dvipsnames,svgnames,x11names}{xcolor}
%
\documentclass[
  letterpaper,
  DIV=11,
  numbers=noendperiod]{scrartcl}
\usepackage{xcolor}
\usepackage{amsmath,amssymb}
\setcounter{secnumdepth}{-\maxdimen} % remove section numbering
\usepackage{iftex}
\ifPDFTeX
  \usepackage[T1]{fontenc}
  \usepackage[utf8]{inputenc}
  \usepackage{textcomp} % provide euro and other symbols
\else % if luatex or xetex
  \usepackage{unicode-math} % this also loads fontspec
  \defaultfontfeatures{Scale=MatchLowercase}
  \defaultfontfeatures[\rmfamily]{Ligatures=TeX,Scale=1}
\fi
\usepackage{lmodern}
\ifPDFTeX\else
  % xetex/luatex font selection
\fi
% Use upquote if available, for straight quotes in verbatim environments
\IfFileExists{upquote.sty}{\usepackage{upquote}}{}
\IfFileExists{microtype.sty}{% use microtype if available
  \usepackage[]{microtype}
  \UseMicrotypeSet[protrusion]{basicmath} % disable protrusion for tt fonts
}{}
\makeatletter
\@ifundefined{KOMAClassName}{% if non-KOMA class
  \IfFileExists{parskip.sty}{%
    \usepackage{parskip}
  }{% else
    \setlength{\parindent}{0pt}
    \setlength{\parskip}{6pt plus 2pt minus 1pt}}
}{% if KOMA class
  \KOMAoptions{parskip=half}}
\makeatother
% Make \paragraph and \subparagraph free-standing
\makeatletter
\ifx\paragraph\undefined\else
  \let\oldparagraph\paragraph
  \renewcommand{\paragraph}{
    \@ifstar
      \xxxParagraphStar
      \xxxParagraphNoStar
  }
  \newcommand{\xxxParagraphStar}[1]{\oldparagraph*{#1}\mbox{}}
  \newcommand{\xxxParagraphNoStar}[1]{\oldparagraph{#1}\mbox{}}
\fi
\ifx\subparagraph\undefined\else
  \let\oldsubparagraph\subparagraph
  \renewcommand{\subparagraph}{
    \@ifstar
      \xxxSubParagraphStar
      \xxxSubParagraphNoStar
  }
  \newcommand{\xxxSubParagraphStar}[1]{\oldsubparagraph*{#1}\mbox{}}
  \newcommand{\xxxSubParagraphNoStar}[1]{\oldsubparagraph{#1}\mbox{}}
\fi
\makeatother


\usepackage{longtable,booktabs,array}
\usepackage{calc} % for calculating minipage widths
% Correct order of tables after \paragraph or \subparagraph
\usepackage{etoolbox}
\makeatletter
\patchcmd\longtable{\par}{\if@noskipsec\mbox{}\fi\par}{}{}
\makeatother
% Allow footnotes in longtable head/foot
\IfFileExists{footnotehyper.sty}{\usepackage{footnotehyper}}{\usepackage{footnote}}
\makesavenoteenv{longtable}
\usepackage{graphicx}
\makeatletter
\newsavebox\pandoc@box
\newcommand*\pandocbounded[1]{% scales image to fit in text height/width
  \sbox\pandoc@box{#1}%
  \Gscale@div\@tempa{\textheight}{\dimexpr\ht\pandoc@box+\dp\pandoc@box\relax}%
  \Gscale@div\@tempb{\linewidth}{\wd\pandoc@box}%
  \ifdim\@tempb\p@<\@tempa\p@\let\@tempa\@tempb\fi% select the smaller of both
  \ifdim\@tempa\p@<\p@\scalebox{\@tempa}{\usebox\pandoc@box}%
  \else\usebox{\pandoc@box}%
  \fi%
}
% Set default figure placement to htbp
\def\fps@figure{htbp}
\makeatother


% definitions for citeproc citations
\NewDocumentCommand\citeproctext{}{}
\NewDocumentCommand\citeproc{mm}{%
  \begingroup\def\citeproctext{#2}\cite{#1}\endgroup}
\makeatletter
 % allow citations to break across lines
 \let\@cite@ofmt\@firstofone
 % avoid brackets around text for \cite:
 \def\@biblabel#1{}
 \def\@cite#1#2{{#1\if@tempswa , #2\fi}}
\makeatother
\newlength{\cslhangindent}
\setlength{\cslhangindent}{1.5em}
\newlength{\csllabelwidth}
\setlength{\csllabelwidth}{3em}
\newenvironment{CSLReferences}[2] % #1 hanging-indent, #2 entry-spacing
 {\begin{list}{}{%
  \setlength{\itemindent}{0pt}
  \setlength{\leftmargin}{0pt}
  \setlength{\parsep}{0pt}
  % turn on hanging indent if param 1 is 1
  \ifodd #1
   \setlength{\leftmargin}{\cslhangindent}
   \setlength{\itemindent}{-1\cslhangindent}
  \fi
  % set entry spacing
  \setlength{\itemsep}{#2\baselineskip}}}
 {\end{list}}
\usepackage{calc}
\newcommand{\CSLBlock}[1]{\hfill\break\parbox[t]{\linewidth}{\strut\ignorespaces#1\strut}}
\newcommand{\CSLLeftMargin}[1]{\parbox[t]{\csllabelwidth}{\strut#1\strut}}
\newcommand{\CSLRightInline}[1]{\parbox[t]{\linewidth - \csllabelwidth}{\strut#1\strut}}
\newcommand{\CSLIndent}[1]{\hspace{\cslhangindent}#1}



\setlength{\emergencystretch}{3em} % prevent overfull lines

\providecommand{\tightlist}{%
  \setlength{\itemsep}{0pt}\setlength{\parskip}{0pt}}



 


\usepackage{parskip}
\usepackage{geometry}
\geometry{margin=1cm}

\usepackage{amsmath}
\usepackage{amssymb}

\usepackage{graphicx}
\setkeys{Gin}{keepaspectratio}
\usepackage{float}
\floatplacement{figure}{H}
\usepackage{caption}
\captionsetup{format=plain,aboveskip=\baselineskip,belowskip=\baselineskip}

\usepackage{xcolor}
\usepackage{hyperref}
\definecolor{urlcolor}{rgb}{0,.145,.698}
\definecolor{linkcolor}{rgb}{.71,0.21,0.01}
\definecolor{citecolor}{rgb}{.12,.54,.11}
\hypersetup{
  breaklinks=true,
  colorlinks=true,
  urlcolor=urlcolor,
  linkcolor=linkcolor,
  citecolor=citecolor
}

\usepackage{orcidlink}
\usepackage{authblk}

\usepackage[bitstream-charter]{mathdesign}
\usepackage[left]{lineno}

% -------------------------
% Shared header: colors + listings + conditional lanes + optional dark mode
% -------------------------

\usepackage{listings}
\usepackage{cite}

% --- Palette ---
\definecolor{bg}{HTML}{111216}      % near-black
\definecolor{fg}{HTML}{FFFFFF}      % white
\definecolor{muted}{HTML}{BBBBBB}
\definecolor{accent}{HTML}{89C2FF}  % link/accent color
\definecolor{accent2}{HTML}{7FE7C4} % second accent

% --- Paths / list bullets ---
\graphicspath{{images/}{./}}
\renewcommand\labelitemi{-}

% -------------------------
% Conditional "sketch lanes"
% -------------------------
% -------------------------
% Conditional "lanes": final / draft / all (AI)
% Usage:
%   \draft{...} \draftNoTag{...} \begin{draftBlock}...\end{draftBlock}
%   \ai{...}    \aiNoTag{...}    \begin{aiBlock}...\end{aiBlock}
%
% Set mode from CLI (examples):
%   pdflatex "\def\docMode{final}% Options for packages loaded elsewhere
% Options for packages loaded elsewhere
\PassOptionsToPackage{unicode}{hyperref}
\PassOptionsToPackage{hyphens}{url}
\PassOptionsToPackage{dvipsnames,svgnames,x11names}{xcolor}
%
\documentclass[
  letterpaper,
  DIV=11,
  numbers=noendperiod]{scrartcl}
\usepackage{xcolor}
\usepackage{amsmath,amssymb}
\setcounter{secnumdepth}{-\maxdimen} % remove section numbering
\usepackage{iftex}
\ifPDFTeX
  \usepackage[T1]{fontenc}
  \usepackage[utf8]{inputenc}
  \usepackage{textcomp} % provide euro and other symbols
\else % if luatex or xetex
  \usepackage{unicode-math} % this also loads fontspec
  \defaultfontfeatures{Scale=MatchLowercase}
  \defaultfontfeatures[\rmfamily]{Ligatures=TeX,Scale=1}
\fi
\usepackage{lmodern}
\ifPDFTeX\else
  % xetex/luatex font selection
\fi
% Use upquote if available, for straight quotes in verbatim environments
\IfFileExists{upquote.sty}{\usepackage{upquote}}{}
\IfFileExists{microtype.sty}{% use microtype if available
  \usepackage[]{microtype}
  \UseMicrotypeSet[protrusion]{basicmath} % disable protrusion for tt fonts
}{}
\makeatletter
\@ifundefined{KOMAClassName}{% if non-KOMA class
  \IfFileExists{parskip.sty}{%
    \usepackage{parskip}
  }{% else
    \setlength{\parindent}{0pt}
    \setlength{\parskip}{6pt plus 2pt minus 1pt}}
}{% if KOMA class
  \KOMAoptions{parskip=half}}
\makeatother
% Make \paragraph and \subparagraph free-standing
\makeatletter
\ifx\paragraph\undefined\else
  \let\oldparagraph\paragraph
  \renewcommand{\paragraph}{
    \@ifstar
      \xxxParagraphStar
      \xxxParagraphNoStar
  }
  \newcommand{\xxxParagraphStar}[1]{\oldparagraph*{#1}\mbox{}}
  \newcommand{\xxxParagraphNoStar}[1]{\oldparagraph{#1}\mbox{}}
\fi
\ifx\subparagraph\undefined\else
  \let\oldsubparagraph\subparagraph
  \renewcommand{\subparagraph}{
    \@ifstar
      \xxxSubParagraphStar
      \xxxSubParagraphNoStar
  }
  \newcommand{\xxxSubParagraphStar}[1]{\oldsubparagraph*{#1}\mbox{}}
  \newcommand{\xxxSubParagraphNoStar}[1]{\oldsubparagraph{#1}\mbox{}}
\fi
\makeatother


\usepackage{longtable,booktabs,array}
\usepackage{calc} % for calculating minipage widths
% Correct order of tables after \paragraph or \subparagraph
\usepackage{etoolbox}
\makeatletter
\patchcmd\longtable{\par}{\if@noskipsec\mbox{}\fi\par}{}{}
\makeatother
% Allow footnotes in longtable head/foot
\IfFileExists{footnotehyper.sty}{\usepackage{footnotehyper}}{\usepackage{footnote}}
\makesavenoteenv{longtable}
\usepackage{graphicx}
\makeatletter
\newsavebox\pandoc@box
\newcommand*\pandocbounded[1]{% scales image to fit in text height/width
  \sbox\pandoc@box{#1}%
  \Gscale@div\@tempa{\textheight}{\dimexpr\ht\pandoc@box+\dp\pandoc@box\relax}%
  \Gscale@div\@tempb{\linewidth}{\wd\pandoc@box}%
  \ifdim\@tempb\p@<\@tempa\p@\let\@tempa\@tempb\fi% select the smaller of both
  \ifdim\@tempa\p@<\p@\scalebox{\@tempa}{\usebox\pandoc@box}%
  \else\usebox{\pandoc@box}%
  \fi%
}
% Set default figure placement to htbp
\def\fps@figure{htbp}
\makeatother


% definitions for citeproc citations
\NewDocumentCommand\citeproctext{}{}
\NewDocumentCommand\citeproc{mm}{%
  \begingroup\def\citeproctext{#2}\cite{#1}\endgroup}
\makeatletter
 % allow citations to break across lines
 \let\@cite@ofmt\@firstofone
 % avoid brackets around text for \cite:
 \def\@biblabel#1{}
 \def\@cite#1#2{{#1\if@tempswa , #2\fi}}
\makeatother
\newlength{\cslhangindent}
\setlength{\cslhangindent}{1.5em}
\newlength{\csllabelwidth}
\setlength{\csllabelwidth}{3em}
\newenvironment{CSLReferences}[2] % #1 hanging-indent, #2 entry-spacing
 {\begin{list}{}{%
  \setlength{\itemindent}{0pt}
  \setlength{\leftmargin}{0pt}
  \setlength{\parsep}{0pt}
  % turn on hanging indent if param 1 is 1
  \ifodd #1
   \setlength{\leftmargin}{\cslhangindent}
   \setlength{\itemindent}{-1\cslhangindent}
  \fi
  % set entry spacing
  \setlength{\itemsep}{#2\baselineskip}}}
 {\end{list}}
\usepackage{calc}
\newcommand{\CSLBlock}[1]{\hfill\break\parbox[t]{\linewidth}{\strut\ignorespaces#1\strut}}
\newcommand{\CSLLeftMargin}[1]{\parbox[t]{\csllabelwidth}{\strut#1\strut}}
\newcommand{\CSLRightInline}[1]{\parbox[t]{\linewidth - \csllabelwidth}{\strut#1\strut}}
\newcommand{\CSLIndent}[1]{\hspace{\cslhangindent}#1}



\setlength{\emergencystretch}{3em} % prevent overfull lines

\providecommand{\tightlist}{%
  \setlength{\itemsep}{0pt}\setlength{\parskip}{0pt}}



 


\usepackage{parskip}
\usepackage{geometry}
\geometry{margin=1cm}

\usepackage{amsmath}
\usepackage{amssymb}

\usepackage{graphicx}
\setkeys{Gin}{keepaspectratio}
\usepackage{float}
\floatplacement{figure}{H}
\usepackage{caption}
\captionsetup{format=plain,aboveskip=\baselineskip,belowskip=\baselineskip}

\usepackage{xcolor}
\usepackage{hyperref}
\definecolor{urlcolor}{rgb}{0,.145,.698}
\definecolor{linkcolor}{rgb}{.71,0.21,0.01}
\definecolor{citecolor}{rgb}{.12,.54,.11}
\hypersetup{
  breaklinks=true,
  colorlinks=true,
  urlcolor=urlcolor,
  linkcolor=linkcolor,
  citecolor=citecolor
}

\usepackage{orcidlink}
\usepackage{authblk}

\usepackage[bitstream-charter]{mathdesign}
\usepackage[left]{lineno}

\input{header.tex}

\providecommand{\tightlist}{\setlength{\itemsep}{0pt}\setlength{\parskip}{0pt}}

\newcommand{\g}{\includegraphics}

\title{Secondary ossification center growth}
\author[1]{Cristian Rodrigo Bustamante-Porras}
\author[1]{Diego Alexander Garzón-Alvarado\,\orcidlink{0000-0003-0072-3738}}
\affil[1]{GNUM Research Group, Department of Mechanical and Mechatronics Engineering, Universidad Nacional de Colombia, Carrera 30 45-03, Bogotá D.C., 111321, Colombia}

\KOMAoption{captions}{tableheading}
\makeatletter
\@ifpackageloaded{caption}{}{\usepackage{caption}}
\AtBeginDocument{%
\ifdefined\contentsname
  \renewcommand*\contentsname{Table of contents}
\else
  \newcommand\contentsname{Table of contents}
\fi
\ifdefined\listfigurename
  \renewcommand*\listfigurename{List of Figures}
\else
  \newcommand\listfigurename{List of Figures}
\fi
\ifdefined\listtablename
  \renewcommand*\listtablename{List of Tables}
\else
  \newcommand\listtablename{List of Tables}
\fi
\ifdefined\figurename
  \renewcommand*\figurename{Figure}
\else
  \newcommand\figurename{Figure}
\fi
\ifdefined\tablename
  \renewcommand*\tablename{Table}
\else
  \newcommand\tablename{Table}
\fi
}
\@ifpackageloaded{float}{}{\usepackage{float}}
\floatstyle{ruled}
\@ifundefined{c@chapter}{\newfloat{codelisting}{h}{lop}}{\newfloat{codelisting}{h}{lop}[chapter]}
\floatname{codelisting}{Listing}
\newcommand*\listoflistings{\listof{codelisting}{List of Listings}}
\makeatother
\makeatletter
\makeatother
\makeatletter
\@ifpackageloaded{caption}{}{\usepackage{caption}}
\@ifpackageloaded{subcaption}{}{\usepackage{subcaption}}
\makeatother
\usepackage{bookmark}
\IfFileExists{xurl.sty}{\usepackage{xurl}}{} % add URL line breaks if available
\urlstyle{same}
\hypersetup{
  pdftitle={Quarto Grid Challenge},
  colorlinks=true,
  linkcolor={blue},
  filecolor={Maroon},
  citecolor={Blue},
  urlcolor={Blue},
  pdfcreator={LaTeX via pandoc}}


\title{Quarto Grid Challenge}
\author{}
\date{}
\begin{document}
\maketitle


\linenumbers

\subsection{Introduction}\label{introduction}

Long bones develop through the process of Endochondral Ossification (EO)
(Kronenberg 2003; Aghajanian and Mohan 2018; Nguyen et al. 2017). This
process starts with the condensation of mesenchymal cells, continuing
with intricate processes of differentiation, proliferation, hypertrophy,
apoptosis and matrix modification (Kronenberg 2003; Aghajanian and Mohan
2018; Nguyen et al. 2017). During this transformation, hyaline cartilage
is replaced by bone, developing primary (POC) and secondary (SOC)
ossification centers; the former allows longitudinal development, while
the latter promotes epiphysis formation. This progression is modulated
by multiple factors, including cellular, biochemical, and mechanical
effects (Carter and Wong 1988; Sundaramurthy and Mao 2006). Disturbances
on this growth mechanism can cause different pathologies in children and
teenagers; therefore, understanding this process is crucial to explain,
prevent and treat these disorders (Comellas and Shefelbine 2022; Ogden,
Hensinger, and McCollough 2000). For this reason, several mathematical
and computational models have been proposed to simulate EO and SOC
emergence. Early works suggested that cyclical mechanical stimuli could
influence ossification, allowing the prediction on location of SOCs
(Carter and Wong 1988).

Additionally, other factors have been considered, e.g., bone remodelling
(Wong and Carter 1990; Stevens, Beaupré, and Carter 1999), biochemical
feedback loops and cellular mechanisms (D. A. Garzón-Alvarado,
García-Aznar, and Doblaré 2009; Diego Alexander Garzón-Alvarado,
Garcı́a-Aznar, and Doblaré 2009; Stoop, Yokoyama, and Adachi 2025), or a
combination of these elements\\
(Peinado Cortés, Vanegas Acosta, and Garzón Alvarado 2011; Vaca-González
et al. 2018; Carrera-Pinzón et al. 2020; Sadeghian, Shapiro, and
Shefelbine 2021). Given that multiple SOCs can appear on a single
epiphysis (Ogden, Hensinger, and McCollough 2000), studies based on
reaction-diffusion modeling have been conducted to analyze the influence
of epiphyseal size on the number of SOCs (Diego Alexander
Garzón-Alvarado 2011). However, previous models have mostly focused on
the emergence of a single SOC, overlooking multiple SOCs presence on the
epiphysis, as observed in the humerus, tibia, femur, and other bones.
Therefore, this study builds on previous work that predicted SOC
locations based primarily on geometric and mechanical factors
(Bustamante-Porras et al. 2025) by incorporating biochemical influences.
It aims to explore how geometry, mechanics, and biochemical conditions
interact to govern SOC formation, including double-center ossification
cases.

\begin{aiBlock}
AIIIIII AAA
\end{aiBlock}
\begin{draftBlock}
better?
\end{draftBlock}

\phantomsection\label{refs}
\begin{CSLReferences}{1}{0}
\bibitem[\citeproctext]{ref-aghajanianArtBuildingBone2018}
Aghajanian, Patrick, and Subburaman Mohan. 2018. {``The Art of Building
Bone: Emerging Role of Chondrocyte-to-Osteoblast Transdifferentiation in
Endochondral Ossification.''} \emph{Bone Research} 6 (1): 19.
\url{https://doi.org/10.1038/s41413-018-0021-z}.

\bibitem[\citeproctext]{ref-bustamante-porrasModelingEndochondralOssification2025}
Bustamante-Porras, Cristian Rodrigo, Kalenia Marquez-Florez, Carlos
Alberto Duque-Daza, and Diego Alexander Garzón-Alvarado. 2025.
{``Modeling Endochondral Ossification: {Effects} of Mechanical Loading
and Bone Shape.''} \emph{Journal of Orthopaedics} 68 (October):
197--218. \url{https://doi.org/10.1016/j.jor.2025.05.034}.

\bibitem[\citeproctext]{ref-carrera-pinzonComputationalModelSynovial2020}
Carrera-Pinzón, Andrés Felipe, Kalenia Márquez-Flórez, Reuben H. Kraft,
Salah Ramtani, and Diego Alexander Garzón-Alvarado. 2020.
{``Computational Model of a Synovial Joint Morphogenesis.''}
\emph{Biomechanics and Modeling in Mechanobiology} 19 (5): 1389--1402.
\url{https://doi.org/10.1007/s10237-019-01277-4}.

\bibitem[\citeproctext]{ref-carterRoleMechanicalLoading1988}
Carter, Dennis R., and Marcy Wong. 1988. {``The Role of Mechanical
Loading Histories in the Development of Diarthrodial Joints.''}
\emph{Journal of Orthopaedic Research} 6 (6): 804--16.
\url{https://doi.org/10.1002/jor.1100060604}.

\bibitem[\citeproctext]{ref-comellasRoleComputationalModels2022}
Comellas, Ester, and Sandra J. Shefelbine. 2022. {``The Role of
Computational Models in Mechanobiology of Growing Bone.''}
\emph{Frontiers in Bioengineering and Biotechnology} 10.
\url{https://doi.org/10.3389/fbioe.2022.973788}.

\bibitem[\citeproctext]{ref-garzon-alvaradoReactiondiffusionModelLong2009}
Garzón-Alvarado, D. A., J. M. García-Aznar, and M. Doblaré. 2009. {``A
Reaction-Diffusion Model for Long Bones Growth.''} \emph{Biomechanics
and Modeling in Mechanobiology} 8 (5): 381--95.
\url{https://doi.org/10.1007/s10237-008-0144-z}.

\bibitem[\citeproctext]{ref-garzon-alvaradoCanSizeEpiphysis2011}
Garzón-Alvarado, Diego Alexander. 2011. {``Can the Size of the Epiphysis
Determine the Number of Secondary Ossification Centers? {A} Mathematical
Approach.''} \emph{Computer Methods in Biomechanics and Biomedical
Engineering} 14 (9): 819--26.

\bibitem[\citeproctext]{ref-garzon-alvaradoAppearanceLocationSecondary2009}
Garzón-Alvarado, Diego Alexander, JM Garcı́a-Aznar, and Manuel Doblaré.
2009. {``Appearance and Location of Secondary Ossification Centres May
Be Explained by a Reaction--Diffusion Mechanism.''} \emph{Computers in
Biology and Medicine} 39 (6): 554--61.

\bibitem[\citeproctext]{ref-kronenbergDevelopmentalRegulationGrowth2003}
Kronenberg, Henry M. 2003. {``Developmental Regulation of the Growth
Plate.''} \emph{Nature} 423 (6937): 332--36.
\url{https://doi.org/10.1038/nature01657}.

\bibitem[\citeproctext]{ref-nguyenImagingPediatricGrowth2017}
Nguyen, Jie C., B. Keegan Markhardt, Arnold C. Merrow, and Jerry R.
Dwek. 2017. {``Imaging of {Pediatric Growth Plate Disturbances}.''}
\emph{RadioGraphics} 37 (6): 1791--1812.
\url{https://doi.org/10.1148/rg.2017170029}.

\bibitem[\citeproctext]{ref-ogdenDiagnosticImaging2000}
Ogden, John A., R. N. Hensinger, and N. McCollough. 2000. {``Diagnostic
{Imaging}.''} In \emph{Skeletal Injury in the Child}, 3rd ed., 128.
Springer.

\bibitem[\citeproctext]{ref-peinadocortesMechanobiologicalModelEpiphysis2011}
Peinado Cortés, L. M., J. C. Vanegas Acosta, and D. A. Garzón Alvarado.
2011. {``A Mechanobiological Model of Epiphysis Structures Formation.''}
\emph{Journal of Theoretical Biology} 287: 13--25.
\url{https://doi.org/10.1016/j.jtbi.2011.07.011}.

\bibitem[\citeproctext]{ref-sadeghianComputationalModelEndochondral2021}
Sadeghian, S. Mahsa, Frederic D. Shapiro, and Sandra J. Shefelbine.
2021. {``Computational Model of Endochondral Ossification: {Simulating}
Growth of a Long Bone.''} \emph{Bone} 153: 116132.
\url{https://doi.org/10.1016/j.bone.2021.116132}.

\bibitem[\citeproctext]{ref-stevensComputerModelEndochondral1999}
Stevens, S. S., G. S. Beaupré, and D. R. Carter. 1999. {``Computer Model
of Endochondral Growth and Ossification in Long Bones: Biological and
Mechanobiological Influences.''} \emph{Journal of Orthopaedic Research :
Official Publication of the Orthopaedic Research Society} 17 (5):
646--53. \url{https://doi.org/10.1002/jor.1100170505}.

\bibitem[\citeproctext]{ref-stoopTimingRestingZone2025}
Stoop, Jorik, Yuka Yokoyama, and Taiji Adachi. 2025. {``Timing of
Resting Zone Parathyroid Hormone-Related Protein Expression Affects
Maintenance of the Growth Plate During Secondary Ossification: A
Computational Study.''} \emph{Biomechanics and Modeling in
Mechanobiology} 24 (1): 125--37.
\url{https://doi.org/10.1007/s10237-024-01899-3}.

\bibitem[\citeproctext]{ref-sundaramurthyModulationEndochondralDevelopment2006}
Sundaramurthy, Sona, and Jeremy J. Mao. 2006. {``Modulation of
Endochondral Development of the Distal Femoral Condyle by Mechanical
Loading.''} \emph{Journal of Orthopaedic Research : Official Publication
of the Orthopaedic Research Society} 24 (2): 229--41.
\url{https://doi.org/10.1002/jor.20024}.

\bibitem[\citeproctext]{ref-vaca-gonzalezMechanobiologicalModelingEndochondral2018}
Vaca-González, Juan Jairo, Miguelangel Moncayo-Donoso, Johana Guevara,
Yoshie Hata, Sandra J. Shefelbine, and Diego Alexander Garzón-Alvarado.
2018. {``Mechanobiological Modeling of Endochondral Ossification: An
Experimental and Computational Analysis.''} \emph{Biomechanics and
Modeling in Mechanobiology} 17 (3): 853--75.
\url{https://doi.org/10.1007/s10237-017-0997-0}.

\bibitem[\citeproctext]{ref-wongTheoreticalModelEndochondral1990}
Wong, M., and D. R. Carter. 1990. {``A Theoretical Model of Endochondral
Ossification and Bone Architectural Construction in Long Bone
Ontogeny.''} \emph{Anatomy and Embryology} 181 (6): 523--32.
\url{https://doi.org/10.1007/BF00174625}.

\end{CSLReferences}




\end{document}
"
%   pdflatex "\def\docMode{draft}% Options for packages loaded elsewhere
% Options for packages loaded elsewhere
\PassOptionsToPackage{unicode}{hyperref}
\PassOptionsToPackage{hyphens}{url}
\PassOptionsToPackage{dvipsnames,svgnames,x11names}{xcolor}
%
\documentclass[
  letterpaper,
  DIV=11,
  numbers=noendperiod]{scrartcl}
\usepackage{xcolor}
\usepackage{amsmath,amssymb}
\setcounter{secnumdepth}{-\maxdimen} % remove section numbering
\usepackage{iftex}
\ifPDFTeX
  \usepackage[T1]{fontenc}
  \usepackage[utf8]{inputenc}
  \usepackage{textcomp} % provide euro and other symbols
\else % if luatex or xetex
  \usepackage{unicode-math} % this also loads fontspec
  \defaultfontfeatures{Scale=MatchLowercase}
  \defaultfontfeatures[\rmfamily]{Ligatures=TeX,Scale=1}
\fi
\usepackage{lmodern}
\ifPDFTeX\else
  % xetex/luatex font selection
\fi
% Use upquote if available, for straight quotes in verbatim environments
\IfFileExists{upquote.sty}{\usepackage{upquote}}{}
\IfFileExists{microtype.sty}{% use microtype if available
  \usepackage[]{microtype}
  \UseMicrotypeSet[protrusion]{basicmath} % disable protrusion for tt fonts
}{}
\makeatletter
\@ifundefined{KOMAClassName}{% if non-KOMA class
  \IfFileExists{parskip.sty}{%
    \usepackage{parskip}
  }{% else
    \setlength{\parindent}{0pt}
    \setlength{\parskip}{6pt plus 2pt minus 1pt}}
}{% if KOMA class
  \KOMAoptions{parskip=half}}
\makeatother
% Make \paragraph and \subparagraph free-standing
\makeatletter
\ifx\paragraph\undefined\else
  \let\oldparagraph\paragraph
  \renewcommand{\paragraph}{
    \@ifstar
      \xxxParagraphStar
      \xxxParagraphNoStar
  }
  \newcommand{\xxxParagraphStar}[1]{\oldparagraph*{#1}\mbox{}}
  \newcommand{\xxxParagraphNoStar}[1]{\oldparagraph{#1}\mbox{}}
\fi
\ifx\subparagraph\undefined\else
  \let\oldsubparagraph\subparagraph
  \renewcommand{\subparagraph}{
    \@ifstar
      \xxxSubParagraphStar
      \xxxSubParagraphNoStar
  }
  \newcommand{\xxxSubParagraphStar}[1]{\oldsubparagraph*{#1}\mbox{}}
  \newcommand{\xxxSubParagraphNoStar}[1]{\oldsubparagraph{#1}\mbox{}}
\fi
\makeatother


\usepackage{longtable,booktabs,array}
\usepackage{calc} % for calculating minipage widths
% Correct order of tables after \paragraph or \subparagraph
\usepackage{etoolbox}
\makeatletter
\patchcmd\longtable{\par}{\if@noskipsec\mbox{}\fi\par}{}{}
\makeatother
% Allow footnotes in longtable head/foot
\IfFileExists{footnotehyper.sty}{\usepackage{footnotehyper}}{\usepackage{footnote}}
\makesavenoteenv{longtable}
\usepackage{graphicx}
\makeatletter
\newsavebox\pandoc@box
\newcommand*\pandocbounded[1]{% scales image to fit in text height/width
  \sbox\pandoc@box{#1}%
  \Gscale@div\@tempa{\textheight}{\dimexpr\ht\pandoc@box+\dp\pandoc@box\relax}%
  \Gscale@div\@tempb{\linewidth}{\wd\pandoc@box}%
  \ifdim\@tempb\p@<\@tempa\p@\let\@tempa\@tempb\fi% select the smaller of both
  \ifdim\@tempa\p@<\p@\scalebox{\@tempa}{\usebox\pandoc@box}%
  \else\usebox{\pandoc@box}%
  \fi%
}
% Set default figure placement to htbp
\def\fps@figure{htbp}
\makeatother


% definitions for citeproc citations
\NewDocumentCommand\citeproctext{}{}
\NewDocumentCommand\citeproc{mm}{%
  \begingroup\def\citeproctext{#2}\cite{#1}\endgroup}
\makeatletter
 % allow citations to break across lines
 \let\@cite@ofmt\@firstofone
 % avoid brackets around text for \cite:
 \def\@biblabel#1{}
 \def\@cite#1#2{{#1\if@tempswa , #2\fi}}
\makeatother
\newlength{\cslhangindent}
\setlength{\cslhangindent}{1.5em}
\newlength{\csllabelwidth}
\setlength{\csllabelwidth}{3em}
\newenvironment{CSLReferences}[2] % #1 hanging-indent, #2 entry-spacing
 {\begin{list}{}{%
  \setlength{\itemindent}{0pt}
  \setlength{\leftmargin}{0pt}
  \setlength{\parsep}{0pt}
  % turn on hanging indent if param 1 is 1
  \ifodd #1
   \setlength{\leftmargin}{\cslhangindent}
   \setlength{\itemindent}{-1\cslhangindent}
  \fi
  % set entry spacing
  \setlength{\itemsep}{#2\baselineskip}}}
 {\end{list}}
\usepackage{calc}
\newcommand{\CSLBlock}[1]{\hfill\break\parbox[t]{\linewidth}{\strut\ignorespaces#1\strut}}
\newcommand{\CSLLeftMargin}[1]{\parbox[t]{\csllabelwidth}{\strut#1\strut}}
\newcommand{\CSLRightInline}[1]{\parbox[t]{\linewidth - \csllabelwidth}{\strut#1\strut}}
\newcommand{\CSLIndent}[1]{\hspace{\cslhangindent}#1}



\setlength{\emergencystretch}{3em} % prevent overfull lines

\providecommand{\tightlist}{%
  \setlength{\itemsep}{0pt}\setlength{\parskip}{0pt}}



 


\usepackage{parskip}
\usepackage{geometry}
\geometry{margin=1cm}

\usepackage{amsmath}
\usepackage{amssymb}

\usepackage{graphicx}
\setkeys{Gin}{keepaspectratio}
\usepackage{float}
\floatplacement{figure}{H}
\usepackage{caption}
\captionsetup{format=plain,aboveskip=\baselineskip,belowskip=\baselineskip}

\usepackage{xcolor}
\usepackage{hyperref}
\definecolor{urlcolor}{rgb}{0,.145,.698}
\definecolor{linkcolor}{rgb}{.71,0.21,0.01}
\definecolor{citecolor}{rgb}{.12,.54,.11}
\hypersetup{
  breaklinks=true,
  colorlinks=true,
  urlcolor=urlcolor,
  linkcolor=linkcolor,
  citecolor=citecolor
}

\usepackage{orcidlink}
\usepackage{authblk}

\usepackage[bitstream-charter]{mathdesign}
\usepackage[left]{lineno}

\input{header.tex}

\providecommand{\tightlist}{\setlength{\itemsep}{0pt}\setlength{\parskip}{0pt}}

\newcommand{\g}{\includegraphics}

\title{Secondary ossification center growth}
\author[1]{Cristian Rodrigo Bustamante-Porras}
\author[1]{Diego Alexander Garzón-Alvarado\,\orcidlink{0000-0003-0072-3738}}
\affil[1]{GNUM Research Group, Department of Mechanical and Mechatronics Engineering, Universidad Nacional de Colombia, Carrera 30 45-03, Bogotá D.C., 111321, Colombia}

\KOMAoption{captions}{tableheading}
\makeatletter
\@ifpackageloaded{caption}{}{\usepackage{caption}}
\AtBeginDocument{%
\ifdefined\contentsname
  \renewcommand*\contentsname{Table of contents}
\else
  \newcommand\contentsname{Table of contents}
\fi
\ifdefined\listfigurename
  \renewcommand*\listfigurename{List of Figures}
\else
  \newcommand\listfigurename{List of Figures}
\fi
\ifdefined\listtablename
  \renewcommand*\listtablename{List of Tables}
\else
  \newcommand\listtablename{List of Tables}
\fi
\ifdefined\figurename
  \renewcommand*\figurename{Figure}
\else
  \newcommand\figurename{Figure}
\fi
\ifdefined\tablename
  \renewcommand*\tablename{Table}
\else
  \newcommand\tablename{Table}
\fi
}
\@ifpackageloaded{float}{}{\usepackage{float}}
\floatstyle{ruled}
\@ifundefined{c@chapter}{\newfloat{codelisting}{h}{lop}}{\newfloat{codelisting}{h}{lop}[chapter]}
\floatname{codelisting}{Listing}
\newcommand*\listoflistings{\listof{codelisting}{List of Listings}}
\makeatother
\makeatletter
\makeatother
\makeatletter
\@ifpackageloaded{caption}{}{\usepackage{caption}}
\@ifpackageloaded{subcaption}{}{\usepackage{subcaption}}
\makeatother
\usepackage{bookmark}
\IfFileExists{xurl.sty}{\usepackage{xurl}}{} % add URL line breaks if available
\urlstyle{same}
\hypersetup{
  pdftitle={Quarto Grid Challenge},
  colorlinks=true,
  linkcolor={blue},
  filecolor={Maroon},
  citecolor={Blue},
  urlcolor={Blue},
  pdfcreator={LaTeX via pandoc}}


\title{Quarto Grid Challenge}
\author{}
\date{}
\begin{document}
\maketitle


\linenumbers

\subsection{Introduction}\label{introduction}

Long bones develop through the process of Endochondral Ossification (EO)
(Kronenberg 2003; Aghajanian and Mohan 2018; Nguyen et al. 2017). This
process starts with the condensation of mesenchymal cells, continuing
with intricate processes of differentiation, proliferation, hypertrophy,
apoptosis and matrix modification (Kronenberg 2003; Aghajanian and Mohan
2018; Nguyen et al. 2017). During this transformation, hyaline cartilage
is replaced by bone, developing primary (POC) and secondary (SOC)
ossification centers; the former allows longitudinal development, while
the latter promotes epiphysis formation. This progression is modulated
by multiple factors, including cellular, biochemical, and mechanical
effects (Carter and Wong 1988; Sundaramurthy and Mao 2006). Disturbances
on this growth mechanism can cause different pathologies in children and
teenagers; therefore, understanding this process is crucial to explain,
prevent and treat these disorders (Comellas and Shefelbine 2022; Ogden,
Hensinger, and McCollough 2000). For this reason, several mathematical
and computational models have been proposed to simulate EO and SOC
emergence. Early works suggested that cyclical mechanical stimuli could
influence ossification, allowing the prediction on location of SOCs
(Carter and Wong 1988).

Additionally, other factors have been considered, e.g., bone remodelling
(Wong and Carter 1990; Stevens, Beaupré, and Carter 1999), biochemical
feedback loops and cellular mechanisms (D. A. Garzón-Alvarado,
García-Aznar, and Doblaré 2009; Diego Alexander Garzón-Alvarado,
Garcı́a-Aznar, and Doblaré 2009; Stoop, Yokoyama, and Adachi 2025), or a
combination of these elements\\
(Peinado Cortés, Vanegas Acosta, and Garzón Alvarado 2011; Vaca-González
et al. 2018; Carrera-Pinzón et al. 2020; Sadeghian, Shapiro, and
Shefelbine 2021). Given that multiple SOCs can appear on a single
epiphysis (Ogden, Hensinger, and McCollough 2000), studies based on
reaction-diffusion modeling have been conducted to analyze the influence
of epiphyseal size on the number of SOCs (Diego Alexander
Garzón-Alvarado 2011). However, previous models have mostly focused on
the emergence of a single SOC, overlooking multiple SOCs presence on the
epiphysis, as observed in the humerus, tibia, femur, and other bones.
Therefore, this study builds on previous work that predicted SOC
locations based primarily on geometric and mechanical factors
(Bustamante-Porras et al. 2025) by incorporating biochemical influences.
It aims to explore how geometry, mechanics, and biochemical conditions
interact to govern SOC formation, including double-center ossification
cases.

\begin{aiBlock}
AIIIIII AAA
\end{aiBlock}
\begin{draftBlock}
better?
\end{draftBlock}

\phantomsection\label{refs}
\begin{CSLReferences}{1}{0}
\bibitem[\citeproctext]{ref-aghajanianArtBuildingBone2018}
Aghajanian, Patrick, and Subburaman Mohan. 2018. {``The Art of Building
Bone: Emerging Role of Chondrocyte-to-Osteoblast Transdifferentiation in
Endochondral Ossification.''} \emph{Bone Research} 6 (1): 19.
\url{https://doi.org/10.1038/s41413-018-0021-z}.

\bibitem[\citeproctext]{ref-bustamante-porrasModelingEndochondralOssification2025}
Bustamante-Porras, Cristian Rodrigo, Kalenia Marquez-Florez, Carlos
Alberto Duque-Daza, and Diego Alexander Garzón-Alvarado. 2025.
{``Modeling Endochondral Ossification: {Effects} of Mechanical Loading
and Bone Shape.''} \emph{Journal of Orthopaedics} 68 (October):
197--218. \url{https://doi.org/10.1016/j.jor.2025.05.034}.

\bibitem[\citeproctext]{ref-carrera-pinzonComputationalModelSynovial2020}
Carrera-Pinzón, Andrés Felipe, Kalenia Márquez-Flórez, Reuben H. Kraft,
Salah Ramtani, and Diego Alexander Garzón-Alvarado. 2020.
{``Computational Model of a Synovial Joint Morphogenesis.''}
\emph{Biomechanics and Modeling in Mechanobiology} 19 (5): 1389--1402.
\url{https://doi.org/10.1007/s10237-019-01277-4}.

\bibitem[\citeproctext]{ref-carterRoleMechanicalLoading1988}
Carter, Dennis R., and Marcy Wong. 1988. {``The Role of Mechanical
Loading Histories in the Development of Diarthrodial Joints.''}
\emph{Journal of Orthopaedic Research} 6 (6): 804--16.
\url{https://doi.org/10.1002/jor.1100060604}.

\bibitem[\citeproctext]{ref-comellasRoleComputationalModels2022}
Comellas, Ester, and Sandra J. Shefelbine. 2022. {``The Role of
Computational Models in Mechanobiology of Growing Bone.''}
\emph{Frontiers in Bioengineering and Biotechnology} 10.
\url{https://doi.org/10.3389/fbioe.2022.973788}.

\bibitem[\citeproctext]{ref-garzon-alvaradoReactiondiffusionModelLong2009}
Garzón-Alvarado, D. A., J. M. García-Aznar, and M. Doblaré. 2009. {``A
Reaction-Diffusion Model for Long Bones Growth.''} \emph{Biomechanics
and Modeling in Mechanobiology} 8 (5): 381--95.
\url{https://doi.org/10.1007/s10237-008-0144-z}.

\bibitem[\citeproctext]{ref-garzon-alvaradoCanSizeEpiphysis2011}
Garzón-Alvarado, Diego Alexander. 2011. {``Can the Size of the Epiphysis
Determine the Number of Secondary Ossification Centers? {A} Mathematical
Approach.''} \emph{Computer Methods in Biomechanics and Biomedical
Engineering} 14 (9): 819--26.

\bibitem[\citeproctext]{ref-garzon-alvaradoAppearanceLocationSecondary2009}
Garzón-Alvarado, Diego Alexander, JM Garcı́a-Aznar, and Manuel Doblaré.
2009. {``Appearance and Location of Secondary Ossification Centres May
Be Explained by a Reaction--Diffusion Mechanism.''} \emph{Computers in
Biology and Medicine} 39 (6): 554--61.

\bibitem[\citeproctext]{ref-kronenbergDevelopmentalRegulationGrowth2003}
Kronenberg, Henry M. 2003. {``Developmental Regulation of the Growth
Plate.''} \emph{Nature} 423 (6937): 332--36.
\url{https://doi.org/10.1038/nature01657}.

\bibitem[\citeproctext]{ref-nguyenImagingPediatricGrowth2017}
Nguyen, Jie C., B. Keegan Markhardt, Arnold C. Merrow, and Jerry R.
Dwek. 2017. {``Imaging of {Pediatric Growth Plate Disturbances}.''}
\emph{RadioGraphics} 37 (6): 1791--1812.
\url{https://doi.org/10.1148/rg.2017170029}.

\bibitem[\citeproctext]{ref-ogdenDiagnosticImaging2000}
Ogden, John A., R. N. Hensinger, and N. McCollough. 2000. {``Diagnostic
{Imaging}.''} In \emph{Skeletal Injury in the Child}, 3rd ed., 128.
Springer.

\bibitem[\citeproctext]{ref-peinadocortesMechanobiologicalModelEpiphysis2011}
Peinado Cortés, L. M., J. C. Vanegas Acosta, and D. A. Garzón Alvarado.
2011. {``A Mechanobiological Model of Epiphysis Structures Formation.''}
\emph{Journal of Theoretical Biology} 287: 13--25.
\url{https://doi.org/10.1016/j.jtbi.2011.07.011}.

\bibitem[\citeproctext]{ref-sadeghianComputationalModelEndochondral2021}
Sadeghian, S. Mahsa, Frederic D. Shapiro, and Sandra J. Shefelbine.
2021. {``Computational Model of Endochondral Ossification: {Simulating}
Growth of a Long Bone.''} \emph{Bone} 153: 116132.
\url{https://doi.org/10.1016/j.bone.2021.116132}.

\bibitem[\citeproctext]{ref-stevensComputerModelEndochondral1999}
Stevens, S. S., G. S. Beaupré, and D. R. Carter. 1999. {``Computer Model
of Endochondral Growth and Ossification in Long Bones: Biological and
Mechanobiological Influences.''} \emph{Journal of Orthopaedic Research :
Official Publication of the Orthopaedic Research Society} 17 (5):
646--53. \url{https://doi.org/10.1002/jor.1100170505}.

\bibitem[\citeproctext]{ref-stoopTimingRestingZone2025}
Stoop, Jorik, Yuka Yokoyama, and Taiji Adachi. 2025. {``Timing of
Resting Zone Parathyroid Hormone-Related Protein Expression Affects
Maintenance of the Growth Plate During Secondary Ossification: A
Computational Study.''} \emph{Biomechanics and Modeling in
Mechanobiology} 24 (1): 125--37.
\url{https://doi.org/10.1007/s10237-024-01899-3}.

\bibitem[\citeproctext]{ref-sundaramurthyModulationEndochondralDevelopment2006}
Sundaramurthy, Sona, and Jeremy J. Mao. 2006. {``Modulation of
Endochondral Development of the Distal Femoral Condyle by Mechanical
Loading.''} \emph{Journal of Orthopaedic Research : Official Publication
of the Orthopaedic Research Society} 24 (2): 229--41.
\url{https://doi.org/10.1002/jor.20024}.

\bibitem[\citeproctext]{ref-vaca-gonzalezMechanobiologicalModelingEndochondral2018}
Vaca-González, Juan Jairo, Miguelangel Moncayo-Donoso, Johana Guevara,
Yoshie Hata, Sandra J. Shefelbine, and Diego Alexander Garzón-Alvarado.
2018. {``Mechanobiological Modeling of Endochondral Ossification: An
Experimental and Computational Analysis.''} \emph{Biomechanics and
Modeling in Mechanobiology} 17 (3): 853--75.
\url{https://doi.org/10.1007/s10237-017-0997-0}.

\bibitem[\citeproctext]{ref-wongTheoreticalModelEndochondral1990}
Wong, M., and D. R. Carter. 1990. {``A Theoretical Model of Endochondral
Ossification and Bone Architectural Construction in Long Bone
Ontogeny.''} \emph{Anatomy and Embryology} 181 (6): 523--32.
\url{https://doi.org/10.1007/BF00174625}.

\end{CSLReferences}




\end{document}
"
%   pdflatex "\def\docMode{all}% Options for packages loaded elsewhere
% Options for packages loaded elsewhere
\PassOptionsToPackage{unicode}{hyperref}
\PassOptionsToPackage{hyphens}{url}
\PassOptionsToPackage{dvipsnames,svgnames,x11names}{xcolor}
%
\documentclass[
  letterpaper,
  DIV=11,
  numbers=noendperiod]{scrartcl}
\usepackage{xcolor}
\usepackage{amsmath,amssymb}
\setcounter{secnumdepth}{-\maxdimen} % remove section numbering
\usepackage{iftex}
\ifPDFTeX
  \usepackage[T1]{fontenc}
  \usepackage[utf8]{inputenc}
  \usepackage{textcomp} % provide euro and other symbols
\else % if luatex or xetex
  \usepackage{unicode-math} % this also loads fontspec
  \defaultfontfeatures{Scale=MatchLowercase}
  \defaultfontfeatures[\rmfamily]{Ligatures=TeX,Scale=1}
\fi
\usepackage{lmodern}
\ifPDFTeX\else
  % xetex/luatex font selection
\fi
% Use upquote if available, for straight quotes in verbatim environments
\IfFileExists{upquote.sty}{\usepackage{upquote}}{}
\IfFileExists{microtype.sty}{% use microtype if available
  \usepackage[]{microtype}
  \UseMicrotypeSet[protrusion]{basicmath} % disable protrusion for tt fonts
}{}
\makeatletter
\@ifundefined{KOMAClassName}{% if non-KOMA class
  \IfFileExists{parskip.sty}{%
    \usepackage{parskip}
  }{% else
    \setlength{\parindent}{0pt}
    \setlength{\parskip}{6pt plus 2pt minus 1pt}}
}{% if KOMA class
  \KOMAoptions{parskip=half}}
\makeatother
% Make \paragraph and \subparagraph free-standing
\makeatletter
\ifx\paragraph\undefined\else
  \let\oldparagraph\paragraph
  \renewcommand{\paragraph}{
    \@ifstar
      \xxxParagraphStar
      \xxxParagraphNoStar
  }
  \newcommand{\xxxParagraphStar}[1]{\oldparagraph*{#1}\mbox{}}
  \newcommand{\xxxParagraphNoStar}[1]{\oldparagraph{#1}\mbox{}}
\fi
\ifx\subparagraph\undefined\else
  \let\oldsubparagraph\subparagraph
  \renewcommand{\subparagraph}{
    \@ifstar
      \xxxSubParagraphStar
      \xxxSubParagraphNoStar
  }
  \newcommand{\xxxSubParagraphStar}[1]{\oldsubparagraph*{#1}\mbox{}}
  \newcommand{\xxxSubParagraphNoStar}[1]{\oldsubparagraph{#1}\mbox{}}
\fi
\makeatother


\usepackage{longtable,booktabs,array}
\usepackage{calc} % for calculating minipage widths
% Correct order of tables after \paragraph or \subparagraph
\usepackage{etoolbox}
\makeatletter
\patchcmd\longtable{\par}{\if@noskipsec\mbox{}\fi\par}{}{}
\makeatother
% Allow footnotes in longtable head/foot
\IfFileExists{footnotehyper.sty}{\usepackage{footnotehyper}}{\usepackage{footnote}}
\makesavenoteenv{longtable}
\usepackage{graphicx}
\makeatletter
\newsavebox\pandoc@box
\newcommand*\pandocbounded[1]{% scales image to fit in text height/width
  \sbox\pandoc@box{#1}%
  \Gscale@div\@tempa{\textheight}{\dimexpr\ht\pandoc@box+\dp\pandoc@box\relax}%
  \Gscale@div\@tempb{\linewidth}{\wd\pandoc@box}%
  \ifdim\@tempb\p@<\@tempa\p@\let\@tempa\@tempb\fi% select the smaller of both
  \ifdim\@tempa\p@<\p@\scalebox{\@tempa}{\usebox\pandoc@box}%
  \else\usebox{\pandoc@box}%
  \fi%
}
% Set default figure placement to htbp
\def\fps@figure{htbp}
\makeatother


% definitions for citeproc citations
\NewDocumentCommand\citeproctext{}{}
\NewDocumentCommand\citeproc{mm}{%
  \begingroup\def\citeproctext{#2}\cite{#1}\endgroup}
\makeatletter
 % allow citations to break across lines
 \let\@cite@ofmt\@firstofone
 % avoid brackets around text for \cite:
 \def\@biblabel#1{}
 \def\@cite#1#2{{#1\if@tempswa , #2\fi}}
\makeatother
\newlength{\cslhangindent}
\setlength{\cslhangindent}{1.5em}
\newlength{\csllabelwidth}
\setlength{\csllabelwidth}{3em}
\newenvironment{CSLReferences}[2] % #1 hanging-indent, #2 entry-spacing
 {\begin{list}{}{%
  \setlength{\itemindent}{0pt}
  \setlength{\leftmargin}{0pt}
  \setlength{\parsep}{0pt}
  % turn on hanging indent if param 1 is 1
  \ifodd #1
   \setlength{\leftmargin}{\cslhangindent}
   \setlength{\itemindent}{-1\cslhangindent}
  \fi
  % set entry spacing
  \setlength{\itemsep}{#2\baselineskip}}}
 {\end{list}}
\usepackage{calc}
\newcommand{\CSLBlock}[1]{\hfill\break\parbox[t]{\linewidth}{\strut\ignorespaces#1\strut}}
\newcommand{\CSLLeftMargin}[1]{\parbox[t]{\csllabelwidth}{\strut#1\strut}}
\newcommand{\CSLRightInline}[1]{\parbox[t]{\linewidth - \csllabelwidth}{\strut#1\strut}}
\newcommand{\CSLIndent}[1]{\hspace{\cslhangindent}#1}



\setlength{\emergencystretch}{3em} % prevent overfull lines

\providecommand{\tightlist}{%
  \setlength{\itemsep}{0pt}\setlength{\parskip}{0pt}}



 


\usepackage{parskip}
\usepackage{geometry}
\geometry{margin=1cm}

\usepackage{amsmath}
\usepackage{amssymb}

\usepackage{graphicx}
\setkeys{Gin}{keepaspectratio}
\usepackage{float}
\floatplacement{figure}{H}
\usepackage{caption}
\captionsetup{format=plain,aboveskip=\baselineskip,belowskip=\baselineskip}

\usepackage{xcolor}
\usepackage{hyperref}
\definecolor{urlcolor}{rgb}{0,.145,.698}
\definecolor{linkcolor}{rgb}{.71,0.21,0.01}
\definecolor{citecolor}{rgb}{.12,.54,.11}
\hypersetup{
  breaklinks=true,
  colorlinks=true,
  urlcolor=urlcolor,
  linkcolor=linkcolor,
  citecolor=citecolor
}

\usepackage{orcidlink}
\usepackage{authblk}

\usepackage[bitstream-charter]{mathdesign}
\usepackage[left]{lineno}

\input{header.tex}

\providecommand{\tightlist}{\setlength{\itemsep}{0pt}\setlength{\parskip}{0pt}}

\newcommand{\g}{\includegraphics}

\title{Secondary ossification center growth}
\author[1]{Cristian Rodrigo Bustamante-Porras}
\author[1]{Diego Alexander Garzón-Alvarado\,\orcidlink{0000-0003-0072-3738}}
\affil[1]{GNUM Research Group, Department of Mechanical and Mechatronics Engineering, Universidad Nacional de Colombia, Carrera 30 45-03, Bogotá D.C., 111321, Colombia}

\KOMAoption{captions}{tableheading}
\makeatletter
\@ifpackageloaded{caption}{}{\usepackage{caption}}
\AtBeginDocument{%
\ifdefined\contentsname
  \renewcommand*\contentsname{Table of contents}
\else
  \newcommand\contentsname{Table of contents}
\fi
\ifdefined\listfigurename
  \renewcommand*\listfigurename{List of Figures}
\else
  \newcommand\listfigurename{List of Figures}
\fi
\ifdefined\listtablename
  \renewcommand*\listtablename{List of Tables}
\else
  \newcommand\listtablename{List of Tables}
\fi
\ifdefined\figurename
  \renewcommand*\figurename{Figure}
\else
  \newcommand\figurename{Figure}
\fi
\ifdefined\tablename
  \renewcommand*\tablename{Table}
\else
  \newcommand\tablename{Table}
\fi
}
\@ifpackageloaded{float}{}{\usepackage{float}}
\floatstyle{ruled}
\@ifundefined{c@chapter}{\newfloat{codelisting}{h}{lop}}{\newfloat{codelisting}{h}{lop}[chapter]}
\floatname{codelisting}{Listing}
\newcommand*\listoflistings{\listof{codelisting}{List of Listings}}
\makeatother
\makeatletter
\makeatother
\makeatletter
\@ifpackageloaded{caption}{}{\usepackage{caption}}
\@ifpackageloaded{subcaption}{}{\usepackage{subcaption}}
\makeatother
\usepackage{bookmark}
\IfFileExists{xurl.sty}{\usepackage{xurl}}{} % add URL line breaks if available
\urlstyle{same}
\hypersetup{
  pdftitle={Quarto Grid Challenge},
  colorlinks=true,
  linkcolor={blue},
  filecolor={Maroon},
  citecolor={Blue},
  urlcolor={Blue},
  pdfcreator={LaTeX via pandoc}}


\title{Quarto Grid Challenge}
\author{}
\date{}
\begin{document}
\maketitle


\linenumbers

\subsection{Introduction}\label{introduction}

Long bones develop through the process of Endochondral Ossification (EO)
(Kronenberg 2003; Aghajanian and Mohan 2018; Nguyen et al. 2017). This
process starts with the condensation of mesenchymal cells, continuing
with intricate processes of differentiation, proliferation, hypertrophy,
apoptosis and matrix modification (Kronenberg 2003; Aghajanian and Mohan
2018; Nguyen et al. 2017). During this transformation, hyaline cartilage
is replaced by bone, developing primary (POC) and secondary (SOC)
ossification centers; the former allows longitudinal development, while
the latter promotes epiphysis formation. This progression is modulated
by multiple factors, including cellular, biochemical, and mechanical
effects (Carter and Wong 1988; Sundaramurthy and Mao 2006). Disturbances
on this growth mechanism can cause different pathologies in children and
teenagers; therefore, understanding this process is crucial to explain,
prevent and treat these disorders (Comellas and Shefelbine 2022; Ogden,
Hensinger, and McCollough 2000). For this reason, several mathematical
and computational models have been proposed to simulate EO and SOC
emergence. Early works suggested that cyclical mechanical stimuli could
influence ossification, allowing the prediction on location of SOCs
(Carter and Wong 1988).

Additionally, other factors have been considered, e.g., bone remodelling
(Wong and Carter 1990; Stevens, Beaupré, and Carter 1999), biochemical
feedback loops and cellular mechanisms (D. A. Garzón-Alvarado,
García-Aznar, and Doblaré 2009; Diego Alexander Garzón-Alvarado,
Garcı́a-Aznar, and Doblaré 2009; Stoop, Yokoyama, and Adachi 2025), or a
combination of these elements\\
(Peinado Cortés, Vanegas Acosta, and Garzón Alvarado 2011; Vaca-González
et al. 2018; Carrera-Pinzón et al. 2020; Sadeghian, Shapiro, and
Shefelbine 2021). Given that multiple SOCs can appear on a single
epiphysis (Ogden, Hensinger, and McCollough 2000), studies based on
reaction-diffusion modeling have been conducted to analyze the influence
of epiphyseal size on the number of SOCs (Diego Alexander
Garzón-Alvarado 2011). However, previous models have mostly focused on
the emergence of a single SOC, overlooking multiple SOCs presence on the
epiphysis, as observed in the humerus, tibia, femur, and other bones.
Therefore, this study builds on previous work that predicted SOC
locations based primarily on geometric and mechanical factors
(Bustamante-Porras et al. 2025) by incorporating biochemical influences.
It aims to explore how geometry, mechanics, and biochemical conditions
interact to govern SOC formation, including double-center ossification
cases.

\begin{aiBlock}
AIIIIII AAA
\end{aiBlock}
\begin{draftBlock}
better?
\end{draftBlock}

\phantomsection\label{refs}
\begin{CSLReferences}{1}{0}
\bibitem[\citeproctext]{ref-aghajanianArtBuildingBone2018}
Aghajanian, Patrick, and Subburaman Mohan. 2018. {``The Art of Building
Bone: Emerging Role of Chondrocyte-to-Osteoblast Transdifferentiation in
Endochondral Ossification.''} \emph{Bone Research} 6 (1): 19.
\url{https://doi.org/10.1038/s41413-018-0021-z}.

\bibitem[\citeproctext]{ref-bustamante-porrasModelingEndochondralOssification2025}
Bustamante-Porras, Cristian Rodrigo, Kalenia Marquez-Florez, Carlos
Alberto Duque-Daza, and Diego Alexander Garzón-Alvarado. 2025.
{``Modeling Endochondral Ossification: {Effects} of Mechanical Loading
and Bone Shape.''} \emph{Journal of Orthopaedics} 68 (October):
197--218. \url{https://doi.org/10.1016/j.jor.2025.05.034}.

\bibitem[\citeproctext]{ref-carrera-pinzonComputationalModelSynovial2020}
Carrera-Pinzón, Andrés Felipe, Kalenia Márquez-Flórez, Reuben H. Kraft,
Salah Ramtani, and Diego Alexander Garzón-Alvarado. 2020.
{``Computational Model of a Synovial Joint Morphogenesis.''}
\emph{Biomechanics and Modeling in Mechanobiology} 19 (5): 1389--1402.
\url{https://doi.org/10.1007/s10237-019-01277-4}.

\bibitem[\citeproctext]{ref-carterRoleMechanicalLoading1988}
Carter, Dennis R., and Marcy Wong. 1988. {``The Role of Mechanical
Loading Histories in the Development of Diarthrodial Joints.''}
\emph{Journal of Orthopaedic Research} 6 (6): 804--16.
\url{https://doi.org/10.1002/jor.1100060604}.

\bibitem[\citeproctext]{ref-comellasRoleComputationalModels2022}
Comellas, Ester, and Sandra J. Shefelbine. 2022. {``The Role of
Computational Models in Mechanobiology of Growing Bone.''}
\emph{Frontiers in Bioengineering and Biotechnology} 10.
\url{https://doi.org/10.3389/fbioe.2022.973788}.

\bibitem[\citeproctext]{ref-garzon-alvaradoReactiondiffusionModelLong2009}
Garzón-Alvarado, D. A., J. M. García-Aznar, and M. Doblaré. 2009. {``A
Reaction-Diffusion Model for Long Bones Growth.''} \emph{Biomechanics
and Modeling in Mechanobiology} 8 (5): 381--95.
\url{https://doi.org/10.1007/s10237-008-0144-z}.

\bibitem[\citeproctext]{ref-garzon-alvaradoCanSizeEpiphysis2011}
Garzón-Alvarado, Diego Alexander. 2011. {``Can the Size of the Epiphysis
Determine the Number of Secondary Ossification Centers? {A} Mathematical
Approach.''} \emph{Computer Methods in Biomechanics and Biomedical
Engineering} 14 (9): 819--26.

\bibitem[\citeproctext]{ref-garzon-alvaradoAppearanceLocationSecondary2009}
Garzón-Alvarado, Diego Alexander, JM Garcı́a-Aznar, and Manuel Doblaré.
2009. {``Appearance and Location of Secondary Ossification Centres May
Be Explained by a Reaction--Diffusion Mechanism.''} \emph{Computers in
Biology and Medicine} 39 (6): 554--61.

\bibitem[\citeproctext]{ref-kronenbergDevelopmentalRegulationGrowth2003}
Kronenberg, Henry M. 2003. {``Developmental Regulation of the Growth
Plate.''} \emph{Nature} 423 (6937): 332--36.
\url{https://doi.org/10.1038/nature01657}.

\bibitem[\citeproctext]{ref-nguyenImagingPediatricGrowth2017}
Nguyen, Jie C., B. Keegan Markhardt, Arnold C. Merrow, and Jerry R.
Dwek. 2017. {``Imaging of {Pediatric Growth Plate Disturbances}.''}
\emph{RadioGraphics} 37 (6): 1791--1812.
\url{https://doi.org/10.1148/rg.2017170029}.

\bibitem[\citeproctext]{ref-ogdenDiagnosticImaging2000}
Ogden, John A., R. N. Hensinger, and N. McCollough. 2000. {``Diagnostic
{Imaging}.''} In \emph{Skeletal Injury in the Child}, 3rd ed., 128.
Springer.

\bibitem[\citeproctext]{ref-peinadocortesMechanobiologicalModelEpiphysis2011}
Peinado Cortés, L. M., J. C. Vanegas Acosta, and D. A. Garzón Alvarado.
2011. {``A Mechanobiological Model of Epiphysis Structures Formation.''}
\emph{Journal of Theoretical Biology} 287: 13--25.
\url{https://doi.org/10.1016/j.jtbi.2011.07.011}.

\bibitem[\citeproctext]{ref-sadeghianComputationalModelEndochondral2021}
Sadeghian, S. Mahsa, Frederic D. Shapiro, and Sandra J. Shefelbine.
2021. {``Computational Model of Endochondral Ossification: {Simulating}
Growth of a Long Bone.''} \emph{Bone} 153: 116132.
\url{https://doi.org/10.1016/j.bone.2021.116132}.

\bibitem[\citeproctext]{ref-stevensComputerModelEndochondral1999}
Stevens, S. S., G. S. Beaupré, and D. R. Carter. 1999. {``Computer Model
of Endochondral Growth and Ossification in Long Bones: Biological and
Mechanobiological Influences.''} \emph{Journal of Orthopaedic Research :
Official Publication of the Orthopaedic Research Society} 17 (5):
646--53. \url{https://doi.org/10.1002/jor.1100170505}.

\bibitem[\citeproctext]{ref-stoopTimingRestingZone2025}
Stoop, Jorik, Yuka Yokoyama, and Taiji Adachi. 2025. {``Timing of
Resting Zone Parathyroid Hormone-Related Protein Expression Affects
Maintenance of the Growth Plate During Secondary Ossification: A
Computational Study.''} \emph{Biomechanics and Modeling in
Mechanobiology} 24 (1): 125--37.
\url{https://doi.org/10.1007/s10237-024-01899-3}.

\bibitem[\citeproctext]{ref-sundaramurthyModulationEndochondralDevelopment2006}
Sundaramurthy, Sona, and Jeremy J. Mao. 2006. {``Modulation of
Endochondral Development of the Distal Femoral Condyle by Mechanical
Loading.''} \emph{Journal of Orthopaedic Research : Official Publication
of the Orthopaedic Research Society} 24 (2): 229--41.
\url{https://doi.org/10.1002/jor.20024}.

\bibitem[\citeproctext]{ref-vaca-gonzalezMechanobiologicalModelingEndochondral2018}
Vaca-González, Juan Jairo, Miguelangel Moncayo-Donoso, Johana Guevara,
Yoshie Hata, Sandra J. Shefelbine, and Diego Alexander Garzón-Alvarado.
2018. {``Mechanobiological Modeling of Endochondral Ossification: An
Experimental and Computational Analysis.''} \emph{Biomechanics and
Modeling in Mechanobiology} 17 (3): 853--75.
\url{https://doi.org/10.1007/s10237-017-0997-0}.

\bibitem[\citeproctext]{ref-wongTheoreticalModelEndochondral1990}
Wong, M., and D. R. Carter. 1990. {``A Theoretical Model of Endochondral
Ossification and Bone Architectural Construction in Long Bone
Ontogeny.''} \emph{Anatomy and Embryology} 181 (6): 523--32.
\url{https://doi.org/10.1007/BF00174625}.

\end{CSLReferences}




\end{document}
"
% -------------------------

\usepackage{etoolbox}
\usepackage{comment}

% Can be overridden from command line:
%   \def\docMode{final} or draft or all
\providecommand{\docMode}{final}

\definecolor{aiColor}{HTML}{1F77B4}
\definecolor{draftColor}{HTML}{D62728}

% Helper: execute #2 if docMode == #1, else execute #3
\newcommand{\ifDocMode}[3]{%
  \ifdefstring{\docMode}{#1}{#2}{#3}%
}

% -------------------------
% Tags (customize once)
% -------------------------
\newcommand{\draftTagText}{DRAFT}
\newcommand{\aiTagText}{AI}

\newcommand{\draftTag}{\textbf{[\draftTagText]}\par}
\newcommand{\aiTag}{\textbf{[\aiTagText]}\par}

% -------------------------
% Inline lanes
% -------------------------
\newcommand{\final}[1]{#1}

\newcommand{\draft}[1]{%
  \ifDocMode{final}{}{%
    \ifDocMode{draft}{\begingroup\color{draftColor}\noindent\draftTag\noindent #1\par\endgroup}{%
      \ifDocMode{all}{\begingroup\color{draftColor}\noindent\draftTag\noindent #1\par\endgroup}{}%
    }%
  }%
}

\newcommand{\draftNoTag}[1]{%
  \ifDocMode{final}{}{%
    \ifDocMode{draft}{\begingroup\color{draftColor}#1\par\endgroup}{%
      \ifDocMode{all}{\begingroup\color{draftColor}#1\par\endgroup}{}%
    }%
  }%
}

\newcommand{\ai}[1]{%
  \ifDocMode{all}{\begingroup\color{aiColor}\noindent\aiTag\noindent #1\par\endgroup}{}%
}

\newcommand{\aiNoTag}[1]{%
  \ifDocMode{all}{\begingroup\color{aiColor}#1\par\endgroup}{}%
}

% -------------------------
% Big block environments
% -------------------------
\ifDocMode{final}{
  \excludecomment{draftBlock}
  \excludecomment{aiBlock}
}{
  \newenvironment{draftBlock}
    {\par\begingroup\color{draftColor}\noindent\draftTag\noindent}
    {\par\endgroup}

  \ifDocMode{draft}{
    \excludecomment{aiBlock}
  }{
    \newenvironment{aiBlock}
      {\par\begingroup\color{aiColor}\noindent\aiTag\noindent}
      {\par\endgroup}
  }
}


% -------------------------
% Dark mode toggle
%   default: off
%   enable via CLI: \def\darkMode{1}
% -------------------------
\providecommand{\darkMode}{0}

% Light-mode default listings (only if dark mode is off)
\lstset{
  backgroundcolor=\color{white},
  basicstyle=\ttfamily\small\color{black},
  keywordstyle=\bfseries\color{blue},
  commentstyle=\itshape\color{gray},
  stringstyle=\color{teal},
  numberstyle=\color{gray},
  rulecolor=\color{gray},
  frame=single,
  framesep=8pt,
  tabsize=2,
  showstringspaces=false
}

% If darkMode == 1, override page + listings
\makeatletter
\@ifundefined{ifdefstring}{%
  % etoolbox should already be loaded by conditionalStyling.tex
}{}
\makeatother
\ifdefstring{\darkMode}{1}{% -------------------------
% Dark mode styling
% Enable from CLI with:
%   pdflatex "\def\darkMode{1}% Options for packages loaded elsewhere
% Options for packages loaded elsewhere
\PassOptionsToPackage{unicode}{hyperref}
\PassOptionsToPackage{hyphens}{url}
\PassOptionsToPackage{dvipsnames,svgnames,x11names}{xcolor}
%
\documentclass[
  letterpaper,
  DIV=11,
  numbers=noendperiod]{scrartcl}
\usepackage{xcolor}
\usepackage{amsmath,amssymb}
\setcounter{secnumdepth}{-\maxdimen} % remove section numbering
\usepackage{iftex}
\ifPDFTeX
  \usepackage[T1]{fontenc}
  \usepackage[utf8]{inputenc}
  \usepackage{textcomp} % provide euro and other symbols
\else % if luatex or xetex
  \usepackage{unicode-math} % this also loads fontspec
  \defaultfontfeatures{Scale=MatchLowercase}
  \defaultfontfeatures[\rmfamily]{Ligatures=TeX,Scale=1}
\fi
\usepackage{lmodern}
\ifPDFTeX\else
  % xetex/luatex font selection
\fi
% Use upquote if available, for straight quotes in verbatim environments
\IfFileExists{upquote.sty}{\usepackage{upquote}}{}
\IfFileExists{microtype.sty}{% use microtype if available
  \usepackage[]{microtype}
  \UseMicrotypeSet[protrusion]{basicmath} % disable protrusion for tt fonts
}{}
\makeatletter
\@ifundefined{KOMAClassName}{% if non-KOMA class
  \IfFileExists{parskip.sty}{%
    \usepackage{parskip}
  }{% else
    \setlength{\parindent}{0pt}
    \setlength{\parskip}{6pt plus 2pt minus 1pt}}
}{% if KOMA class
  \KOMAoptions{parskip=half}}
\makeatother
% Make \paragraph and \subparagraph free-standing
\makeatletter
\ifx\paragraph\undefined\else
  \let\oldparagraph\paragraph
  \renewcommand{\paragraph}{
    \@ifstar
      \xxxParagraphStar
      \xxxParagraphNoStar
  }
  \newcommand{\xxxParagraphStar}[1]{\oldparagraph*{#1}\mbox{}}
  \newcommand{\xxxParagraphNoStar}[1]{\oldparagraph{#1}\mbox{}}
\fi
\ifx\subparagraph\undefined\else
  \let\oldsubparagraph\subparagraph
  \renewcommand{\subparagraph}{
    \@ifstar
      \xxxSubParagraphStar
      \xxxSubParagraphNoStar
  }
  \newcommand{\xxxSubParagraphStar}[1]{\oldsubparagraph*{#1}\mbox{}}
  \newcommand{\xxxSubParagraphNoStar}[1]{\oldsubparagraph{#1}\mbox{}}
\fi
\makeatother


\usepackage{longtable,booktabs,array}
\usepackage{calc} % for calculating minipage widths
% Correct order of tables after \paragraph or \subparagraph
\usepackage{etoolbox}
\makeatletter
\patchcmd\longtable{\par}{\if@noskipsec\mbox{}\fi\par}{}{}
\makeatother
% Allow footnotes in longtable head/foot
\IfFileExists{footnotehyper.sty}{\usepackage{footnotehyper}}{\usepackage{footnote}}
\makesavenoteenv{longtable}
\usepackage{graphicx}
\makeatletter
\newsavebox\pandoc@box
\newcommand*\pandocbounded[1]{% scales image to fit in text height/width
  \sbox\pandoc@box{#1}%
  \Gscale@div\@tempa{\textheight}{\dimexpr\ht\pandoc@box+\dp\pandoc@box\relax}%
  \Gscale@div\@tempb{\linewidth}{\wd\pandoc@box}%
  \ifdim\@tempb\p@<\@tempa\p@\let\@tempa\@tempb\fi% select the smaller of both
  \ifdim\@tempa\p@<\p@\scalebox{\@tempa}{\usebox\pandoc@box}%
  \else\usebox{\pandoc@box}%
  \fi%
}
% Set default figure placement to htbp
\def\fps@figure{htbp}
\makeatother


% definitions for citeproc citations
\NewDocumentCommand\citeproctext{}{}
\NewDocumentCommand\citeproc{mm}{%
  \begingroup\def\citeproctext{#2}\cite{#1}\endgroup}
\makeatletter
 % allow citations to break across lines
 \let\@cite@ofmt\@firstofone
 % avoid brackets around text for \cite:
 \def\@biblabel#1{}
 \def\@cite#1#2{{#1\if@tempswa , #2\fi}}
\makeatother
\newlength{\cslhangindent}
\setlength{\cslhangindent}{1.5em}
\newlength{\csllabelwidth}
\setlength{\csllabelwidth}{3em}
\newenvironment{CSLReferences}[2] % #1 hanging-indent, #2 entry-spacing
 {\begin{list}{}{%
  \setlength{\itemindent}{0pt}
  \setlength{\leftmargin}{0pt}
  \setlength{\parsep}{0pt}
  % turn on hanging indent if param 1 is 1
  \ifodd #1
   \setlength{\leftmargin}{\cslhangindent}
   \setlength{\itemindent}{-1\cslhangindent}
  \fi
  % set entry spacing
  \setlength{\itemsep}{#2\baselineskip}}}
 {\end{list}}
\usepackage{calc}
\newcommand{\CSLBlock}[1]{\hfill\break\parbox[t]{\linewidth}{\strut\ignorespaces#1\strut}}
\newcommand{\CSLLeftMargin}[1]{\parbox[t]{\csllabelwidth}{\strut#1\strut}}
\newcommand{\CSLRightInline}[1]{\parbox[t]{\linewidth - \csllabelwidth}{\strut#1\strut}}
\newcommand{\CSLIndent}[1]{\hspace{\cslhangindent}#1}



\setlength{\emergencystretch}{3em} % prevent overfull lines

\providecommand{\tightlist}{%
  \setlength{\itemsep}{0pt}\setlength{\parskip}{0pt}}



 


\usepackage{parskip}
\usepackage{geometry}
\geometry{margin=1cm}

\usepackage{amsmath}
\usepackage{amssymb}

\usepackage{graphicx}
\setkeys{Gin}{keepaspectratio}
\usepackage{float}
\floatplacement{figure}{H}
\usepackage{caption}
\captionsetup{format=plain,aboveskip=\baselineskip,belowskip=\baselineskip}

\usepackage{xcolor}
\usepackage{hyperref}
\definecolor{urlcolor}{rgb}{0,.145,.698}
\definecolor{linkcolor}{rgb}{.71,0.21,0.01}
\definecolor{citecolor}{rgb}{.12,.54,.11}
\hypersetup{
  breaklinks=true,
  colorlinks=true,
  urlcolor=urlcolor,
  linkcolor=linkcolor,
  citecolor=citecolor
}

\usepackage{orcidlink}
\usepackage{authblk}

\usepackage[bitstream-charter]{mathdesign}
\usepackage[left]{lineno}

\input{header.tex}

\providecommand{\tightlist}{\setlength{\itemsep}{0pt}\setlength{\parskip}{0pt}}

\newcommand{\g}{\includegraphics}

\title{Secondary ossification center growth}
\author[1]{Cristian Rodrigo Bustamante-Porras}
\author[1]{Diego Alexander Garzón-Alvarado\,\orcidlink{0000-0003-0072-3738}}
\affil[1]{GNUM Research Group, Department of Mechanical and Mechatronics Engineering, Universidad Nacional de Colombia, Carrera 30 45-03, Bogotá D.C., 111321, Colombia}

\KOMAoption{captions}{tableheading}
\makeatletter
\@ifpackageloaded{caption}{}{\usepackage{caption}}
\AtBeginDocument{%
\ifdefined\contentsname
  \renewcommand*\contentsname{Table of contents}
\else
  \newcommand\contentsname{Table of contents}
\fi
\ifdefined\listfigurename
  \renewcommand*\listfigurename{List of Figures}
\else
  \newcommand\listfigurename{List of Figures}
\fi
\ifdefined\listtablename
  \renewcommand*\listtablename{List of Tables}
\else
  \newcommand\listtablename{List of Tables}
\fi
\ifdefined\figurename
  \renewcommand*\figurename{Figure}
\else
  \newcommand\figurename{Figure}
\fi
\ifdefined\tablename
  \renewcommand*\tablename{Table}
\else
  \newcommand\tablename{Table}
\fi
}
\@ifpackageloaded{float}{}{\usepackage{float}}
\floatstyle{ruled}
\@ifundefined{c@chapter}{\newfloat{codelisting}{h}{lop}}{\newfloat{codelisting}{h}{lop}[chapter]}
\floatname{codelisting}{Listing}
\newcommand*\listoflistings{\listof{codelisting}{List of Listings}}
\makeatother
\makeatletter
\makeatother
\makeatletter
\@ifpackageloaded{caption}{}{\usepackage{caption}}
\@ifpackageloaded{subcaption}{}{\usepackage{subcaption}}
\makeatother
\usepackage{bookmark}
\IfFileExists{xurl.sty}{\usepackage{xurl}}{} % add URL line breaks if available
\urlstyle{same}
\hypersetup{
  pdftitle={Quarto Grid Challenge},
  colorlinks=true,
  linkcolor={blue},
  filecolor={Maroon},
  citecolor={Blue},
  urlcolor={Blue},
  pdfcreator={LaTeX via pandoc}}


\title{Quarto Grid Challenge}
\author{}
\date{}
\begin{document}
\maketitle


\linenumbers

\subsection{Introduction}\label{introduction}

Long bones develop through the process of Endochondral Ossification (EO)
(Kronenberg 2003; Aghajanian and Mohan 2018; Nguyen et al. 2017). This
process starts with the condensation of mesenchymal cells, continuing
with intricate processes of differentiation, proliferation, hypertrophy,
apoptosis and matrix modification (Kronenberg 2003; Aghajanian and Mohan
2018; Nguyen et al. 2017). During this transformation, hyaline cartilage
is replaced by bone, developing primary (POC) and secondary (SOC)
ossification centers; the former allows longitudinal development, while
the latter promotes epiphysis formation. This progression is modulated
by multiple factors, including cellular, biochemical, and mechanical
effects (Carter and Wong 1988; Sundaramurthy and Mao 2006). Disturbances
on this growth mechanism can cause different pathologies in children and
teenagers; therefore, understanding this process is crucial to explain,
prevent and treat these disorders (Comellas and Shefelbine 2022; Ogden,
Hensinger, and McCollough 2000). For this reason, several mathematical
and computational models have been proposed to simulate EO and SOC
emergence. Early works suggested that cyclical mechanical stimuli could
influence ossification, allowing the prediction on location of SOCs
(Carter and Wong 1988).

Additionally, other factors have been considered, e.g., bone remodelling
(Wong and Carter 1990; Stevens, Beaupré, and Carter 1999), biochemical
feedback loops and cellular mechanisms (D. A. Garzón-Alvarado,
García-Aznar, and Doblaré 2009; Diego Alexander Garzón-Alvarado,
Garcı́a-Aznar, and Doblaré 2009; Stoop, Yokoyama, and Adachi 2025), or a
combination of these elements\\
(Peinado Cortés, Vanegas Acosta, and Garzón Alvarado 2011; Vaca-González
et al. 2018; Carrera-Pinzón et al. 2020; Sadeghian, Shapiro, and
Shefelbine 2021). Given that multiple SOCs can appear on a single
epiphysis (Ogden, Hensinger, and McCollough 2000), studies based on
reaction-diffusion modeling have been conducted to analyze the influence
of epiphyseal size on the number of SOCs (Diego Alexander
Garzón-Alvarado 2011). However, previous models have mostly focused on
the emergence of a single SOC, overlooking multiple SOCs presence on the
epiphysis, as observed in the humerus, tibia, femur, and other bones.
Therefore, this study builds on previous work that predicted SOC
locations based primarily on geometric and mechanical factors
(Bustamante-Porras et al. 2025) by incorporating biochemical influences.
It aims to explore how geometry, mechanics, and biochemical conditions
interact to govern SOC formation, including double-center ossification
cases.

\begin{aiBlock}
AIIIIII AAA
\end{aiBlock}
\begin{draftBlock}
better?
\end{draftBlock}

\phantomsection\label{refs}
\begin{CSLReferences}{1}{0}
\bibitem[\citeproctext]{ref-aghajanianArtBuildingBone2018}
Aghajanian, Patrick, and Subburaman Mohan. 2018. {``The Art of Building
Bone: Emerging Role of Chondrocyte-to-Osteoblast Transdifferentiation in
Endochondral Ossification.''} \emph{Bone Research} 6 (1): 19.
\url{https://doi.org/10.1038/s41413-018-0021-z}.

\bibitem[\citeproctext]{ref-bustamante-porrasModelingEndochondralOssification2025}
Bustamante-Porras, Cristian Rodrigo, Kalenia Marquez-Florez, Carlos
Alberto Duque-Daza, and Diego Alexander Garzón-Alvarado. 2025.
{``Modeling Endochondral Ossification: {Effects} of Mechanical Loading
and Bone Shape.''} \emph{Journal of Orthopaedics} 68 (October):
197--218. \url{https://doi.org/10.1016/j.jor.2025.05.034}.

\bibitem[\citeproctext]{ref-carrera-pinzonComputationalModelSynovial2020}
Carrera-Pinzón, Andrés Felipe, Kalenia Márquez-Flórez, Reuben H. Kraft,
Salah Ramtani, and Diego Alexander Garzón-Alvarado. 2020.
{``Computational Model of a Synovial Joint Morphogenesis.''}
\emph{Biomechanics and Modeling in Mechanobiology} 19 (5): 1389--1402.
\url{https://doi.org/10.1007/s10237-019-01277-4}.

\bibitem[\citeproctext]{ref-carterRoleMechanicalLoading1988}
Carter, Dennis R., and Marcy Wong. 1988. {``The Role of Mechanical
Loading Histories in the Development of Diarthrodial Joints.''}
\emph{Journal of Orthopaedic Research} 6 (6): 804--16.
\url{https://doi.org/10.1002/jor.1100060604}.

\bibitem[\citeproctext]{ref-comellasRoleComputationalModels2022}
Comellas, Ester, and Sandra J. Shefelbine. 2022. {``The Role of
Computational Models in Mechanobiology of Growing Bone.''}
\emph{Frontiers in Bioengineering and Biotechnology} 10.
\url{https://doi.org/10.3389/fbioe.2022.973788}.

\bibitem[\citeproctext]{ref-garzon-alvaradoReactiondiffusionModelLong2009}
Garzón-Alvarado, D. A., J. M. García-Aznar, and M. Doblaré. 2009. {``A
Reaction-Diffusion Model for Long Bones Growth.''} \emph{Biomechanics
and Modeling in Mechanobiology} 8 (5): 381--95.
\url{https://doi.org/10.1007/s10237-008-0144-z}.

\bibitem[\citeproctext]{ref-garzon-alvaradoCanSizeEpiphysis2011}
Garzón-Alvarado, Diego Alexander. 2011. {``Can the Size of the Epiphysis
Determine the Number of Secondary Ossification Centers? {A} Mathematical
Approach.''} \emph{Computer Methods in Biomechanics and Biomedical
Engineering} 14 (9): 819--26.

\bibitem[\citeproctext]{ref-garzon-alvaradoAppearanceLocationSecondary2009}
Garzón-Alvarado, Diego Alexander, JM Garcı́a-Aznar, and Manuel Doblaré.
2009. {``Appearance and Location of Secondary Ossification Centres May
Be Explained by a Reaction--Diffusion Mechanism.''} \emph{Computers in
Biology and Medicine} 39 (6): 554--61.

\bibitem[\citeproctext]{ref-kronenbergDevelopmentalRegulationGrowth2003}
Kronenberg, Henry M. 2003. {``Developmental Regulation of the Growth
Plate.''} \emph{Nature} 423 (6937): 332--36.
\url{https://doi.org/10.1038/nature01657}.

\bibitem[\citeproctext]{ref-nguyenImagingPediatricGrowth2017}
Nguyen, Jie C., B. Keegan Markhardt, Arnold C. Merrow, and Jerry R.
Dwek. 2017. {``Imaging of {Pediatric Growth Plate Disturbances}.''}
\emph{RadioGraphics} 37 (6): 1791--1812.
\url{https://doi.org/10.1148/rg.2017170029}.

\bibitem[\citeproctext]{ref-ogdenDiagnosticImaging2000}
Ogden, John A., R. N. Hensinger, and N. McCollough. 2000. {``Diagnostic
{Imaging}.''} In \emph{Skeletal Injury in the Child}, 3rd ed., 128.
Springer.

\bibitem[\citeproctext]{ref-peinadocortesMechanobiologicalModelEpiphysis2011}
Peinado Cortés, L. M., J. C. Vanegas Acosta, and D. A. Garzón Alvarado.
2011. {``A Mechanobiological Model of Epiphysis Structures Formation.''}
\emph{Journal of Theoretical Biology} 287: 13--25.
\url{https://doi.org/10.1016/j.jtbi.2011.07.011}.

\bibitem[\citeproctext]{ref-sadeghianComputationalModelEndochondral2021}
Sadeghian, S. Mahsa, Frederic D. Shapiro, and Sandra J. Shefelbine.
2021. {``Computational Model of Endochondral Ossification: {Simulating}
Growth of a Long Bone.''} \emph{Bone} 153: 116132.
\url{https://doi.org/10.1016/j.bone.2021.116132}.

\bibitem[\citeproctext]{ref-stevensComputerModelEndochondral1999}
Stevens, S. S., G. S. Beaupré, and D. R. Carter. 1999. {``Computer Model
of Endochondral Growth and Ossification in Long Bones: Biological and
Mechanobiological Influences.''} \emph{Journal of Orthopaedic Research :
Official Publication of the Orthopaedic Research Society} 17 (5):
646--53. \url{https://doi.org/10.1002/jor.1100170505}.

\bibitem[\citeproctext]{ref-stoopTimingRestingZone2025}
Stoop, Jorik, Yuka Yokoyama, and Taiji Adachi. 2025. {``Timing of
Resting Zone Parathyroid Hormone-Related Protein Expression Affects
Maintenance of the Growth Plate During Secondary Ossification: A
Computational Study.''} \emph{Biomechanics and Modeling in
Mechanobiology} 24 (1): 125--37.
\url{https://doi.org/10.1007/s10237-024-01899-3}.

\bibitem[\citeproctext]{ref-sundaramurthyModulationEndochondralDevelopment2006}
Sundaramurthy, Sona, and Jeremy J. Mao. 2006. {``Modulation of
Endochondral Development of the Distal Femoral Condyle by Mechanical
Loading.''} \emph{Journal of Orthopaedic Research : Official Publication
of the Orthopaedic Research Society} 24 (2): 229--41.
\url{https://doi.org/10.1002/jor.20024}.

\bibitem[\citeproctext]{ref-vaca-gonzalezMechanobiologicalModelingEndochondral2018}
Vaca-González, Juan Jairo, Miguelangel Moncayo-Donoso, Johana Guevara,
Yoshie Hata, Sandra J. Shefelbine, and Diego Alexander Garzón-Alvarado.
2018. {``Mechanobiological Modeling of Endochondral Ossification: An
Experimental and Computational Analysis.''} \emph{Biomechanics and
Modeling in Mechanobiology} 17 (3): 853--75.
\url{https://doi.org/10.1007/s10237-017-0997-0}.

\bibitem[\citeproctext]{ref-wongTheoreticalModelEndochondral1990}
Wong, M., and D. R. Carter. 1990. {``A Theoretical Model of Endochondral
Ossification and Bone Architectural Construction in Long Bone
Ontogeny.''} \emph{Anatomy and Embryology} 181 (6): 523--32.
\url{https://doi.org/10.1007/BF00174625}.

\end{CSLReferences}




\end{document}
"
% -------------------------

% Softer draft color for dark mode (light brown / tan)
\definecolor{draftColor}{HTML}{C8A06A}

% Background override for dark theme
\definecolor{bg}{HTML}{212121}

\pagecolor{bg}
\color{fg}

\lstset{
  backgroundcolor=\color{bg},
  basicstyle=\ttfamily\small\color{fg},
  keywordstyle=\bfseries\color{accent},
  commentstyle=\itshape\color{muted},
  stringstyle=\color{accent2},
  numberstyle=\color{muted},
  rulecolor=\color{muted},
  frame=single,
  framesep=8pt,
  tabsize=2,
  showstringspaces=false
}
}{}


\providecommand{\tightlist}{\setlength{\itemsep}{0pt}\setlength{\parskip}{0pt}}

\newcommand{\g}{\includegraphics}

\title{Secondary ossification center growth}
\author[1]{Cristian Rodrigo Bustamante-Porras}
\author[1]{Diego Alexander Garzón-Alvarado\,\orcidlink{0000-0003-0072-3738}}
\affil[1]{GNUM Research Group, Department of Mechanical and Mechatronics Engineering, Universidad Nacional de Colombia, Carrera 30 45-03, Bogotá D.C., 111321, Colombia}

\KOMAoption{captions}{tableheading}
\makeatletter
\@ifpackageloaded{caption}{}{\usepackage{caption}}
\AtBeginDocument{%
\ifdefined\contentsname
  \renewcommand*\contentsname{Table of contents}
\else
  \newcommand\contentsname{Table of contents}
\fi
\ifdefined\listfigurename
  \renewcommand*\listfigurename{List of Figures}
\else
  \newcommand\listfigurename{List of Figures}
\fi
\ifdefined\listtablename
  \renewcommand*\listtablename{List of Tables}
\else
  \newcommand\listtablename{List of Tables}
\fi
\ifdefined\figurename
  \renewcommand*\figurename{Figure}
\else
  \newcommand\figurename{Figure}
\fi
\ifdefined\tablename
  \renewcommand*\tablename{Table}
\else
  \newcommand\tablename{Table}
\fi
}
\@ifpackageloaded{float}{}{\usepackage{float}}
\floatstyle{ruled}
\@ifundefined{c@chapter}{\newfloat{codelisting}{h}{lop}}{\newfloat{codelisting}{h}{lop}[chapter]}
\floatname{codelisting}{Listing}
\newcommand*\listoflistings{\listof{codelisting}{List of Listings}}
\makeatother
\makeatletter
\makeatother
\makeatletter
\@ifpackageloaded{caption}{}{\usepackage{caption}}
\@ifpackageloaded{subcaption}{}{\usepackage{subcaption}}
\makeatother
\usepackage{bookmark}
\IfFileExists{xurl.sty}{\usepackage{xurl}}{} % add URL line breaks if available
\urlstyle{same}
\hypersetup{
  pdftitle={Quarto Grid Challenge},
  colorlinks=true,
  linkcolor={blue},
  filecolor={Maroon},
  citecolor={Blue},
  urlcolor={Blue},
  pdfcreator={LaTeX via pandoc}}


\title{Quarto Grid Challenge}
\author{}
\date{}
\begin{document}
\maketitle


\linenumbers

\subsection{Introduction}\label{introduction}

Long bones develop through the process of Endochondral Ossification (EO)
(Kronenberg 2003; Aghajanian and Mohan 2018; Nguyen et al. 2017). This
process starts with the condensation of mesenchymal cells, continuing
with intricate processes of differentiation, proliferation, hypertrophy,
apoptosis and matrix modification (Kronenberg 2003; Aghajanian and Mohan
2018; Nguyen et al. 2017). During this transformation, hyaline cartilage
is replaced by bone, developing primary (POC) and secondary (SOC)
ossification centers; the former allows longitudinal development, while
the latter promotes epiphysis formation. This progression is modulated
by multiple factors, including cellular, biochemical, and mechanical
effects (Carter and Wong 1988; Sundaramurthy and Mao 2006). Disturbances
on this growth mechanism can cause different pathologies in children and
teenagers; therefore, understanding this process is crucial to explain,
prevent and treat these disorders (Comellas and Shefelbine 2022; Ogden,
Hensinger, and McCollough 2000). For this reason, several mathematical
and computational models have been proposed to simulate EO and SOC
emergence. Early works suggested that cyclical mechanical stimuli could
influence ossification, allowing the prediction on location of SOCs
(Carter and Wong 1988).

Additionally, other factors have been considered, e.g., bone remodelling
(Wong and Carter 1990; Stevens, Beaupré, and Carter 1999), biochemical
feedback loops and cellular mechanisms (D. A. Garzón-Alvarado,
García-Aznar, and Doblaré 2009; Diego Alexander Garzón-Alvarado,
Garcı́a-Aznar, and Doblaré 2009; Stoop, Yokoyama, and Adachi 2025), or a
combination of these elements\\
(Peinado Cortés, Vanegas Acosta, and Garzón Alvarado 2011; Vaca-González
et al. 2018; Carrera-Pinzón et al. 2020; Sadeghian, Shapiro, and
Shefelbine 2021). Given that multiple SOCs can appear on a single
epiphysis (Ogden, Hensinger, and McCollough 2000), studies based on
reaction-diffusion modeling have been conducted to analyze the influence
of epiphyseal size on the number of SOCs (Diego Alexander
Garzón-Alvarado 2011). However, previous models have mostly focused on
the emergence of a single SOC, overlooking multiple SOCs presence on the
epiphysis, as observed in the humerus, tibia, femur, and other bones.
Therefore, this study builds on previous work that predicted SOC
locations based primarily on geometric and mechanical factors
(Bustamante-Porras et al. 2025) by incorporating biochemical influences.
It aims to explore how geometry, mechanics, and biochemical conditions
interact to govern SOC formation, including double-center ossification
cases.

\begin{aiBlock}
AIIIIII AAA
\end{aiBlock}
\begin{draftBlock}
better?
\end{draftBlock}

\phantomsection\label{refs}
\begin{CSLReferences}{1}{0}
\bibitem[\citeproctext]{ref-aghajanianArtBuildingBone2018}
Aghajanian, Patrick, and Subburaman Mohan. 2018. {``The Art of Building
Bone: Emerging Role of Chondrocyte-to-Osteoblast Transdifferentiation in
Endochondral Ossification.''} \emph{Bone Research} 6 (1): 19.
\url{https://doi.org/10.1038/s41413-018-0021-z}.

\bibitem[\citeproctext]{ref-bustamante-porrasModelingEndochondralOssification2025}
Bustamante-Porras, Cristian Rodrigo, Kalenia Marquez-Florez, Carlos
Alberto Duque-Daza, and Diego Alexander Garzón-Alvarado. 2025.
{``Modeling Endochondral Ossification: {Effects} of Mechanical Loading
and Bone Shape.''} \emph{Journal of Orthopaedics} 68 (October):
197--218. \url{https://doi.org/10.1016/j.jor.2025.05.034}.

\bibitem[\citeproctext]{ref-carrera-pinzonComputationalModelSynovial2020}
Carrera-Pinzón, Andrés Felipe, Kalenia Márquez-Flórez, Reuben H. Kraft,
Salah Ramtani, and Diego Alexander Garzón-Alvarado. 2020.
{``Computational Model of a Synovial Joint Morphogenesis.''}
\emph{Biomechanics and Modeling in Mechanobiology} 19 (5): 1389--1402.
\url{https://doi.org/10.1007/s10237-019-01277-4}.

\bibitem[\citeproctext]{ref-carterRoleMechanicalLoading1988}
Carter, Dennis R., and Marcy Wong. 1988. {``The Role of Mechanical
Loading Histories in the Development of Diarthrodial Joints.''}
\emph{Journal of Orthopaedic Research} 6 (6): 804--16.
\url{https://doi.org/10.1002/jor.1100060604}.

\bibitem[\citeproctext]{ref-comellasRoleComputationalModels2022}
Comellas, Ester, and Sandra J. Shefelbine. 2022. {``The Role of
Computational Models in Mechanobiology of Growing Bone.''}
\emph{Frontiers in Bioengineering and Biotechnology} 10.
\url{https://doi.org/10.3389/fbioe.2022.973788}.

\bibitem[\citeproctext]{ref-garzon-alvaradoReactiondiffusionModelLong2009}
Garzón-Alvarado, D. A., J. M. García-Aznar, and M. Doblaré. 2009. {``A
Reaction-Diffusion Model for Long Bones Growth.''} \emph{Biomechanics
and Modeling in Mechanobiology} 8 (5): 381--95.
\url{https://doi.org/10.1007/s10237-008-0144-z}.

\bibitem[\citeproctext]{ref-garzon-alvaradoCanSizeEpiphysis2011}
Garzón-Alvarado, Diego Alexander. 2011. {``Can the Size of the Epiphysis
Determine the Number of Secondary Ossification Centers? {A} Mathematical
Approach.''} \emph{Computer Methods in Biomechanics and Biomedical
Engineering} 14 (9): 819--26.

\bibitem[\citeproctext]{ref-garzon-alvaradoAppearanceLocationSecondary2009}
Garzón-Alvarado, Diego Alexander, JM Garcı́a-Aznar, and Manuel Doblaré.
2009. {``Appearance and Location of Secondary Ossification Centres May
Be Explained by a Reaction--Diffusion Mechanism.''} \emph{Computers in
Biology and Medicine} 39 (6): 554--61.

\bibitem[\citeproctext]{ref-kronenbergDevelopmentalRegulationGrowth2003}
Kronenberg, Henry M. 2003. {``Developmental Regulation of the Growth
Plate.''} \emph{Nature} 423 (6937): 332--36.
\url{https://doi.org/10.1038/nature01657}.

\bibitem[\citeproctext]{ref-nguyenImagingPediatricGrowth2017}
Nguyen, Jie C., B. Keegan Markhardt, Arnold C. Merrow, and Jerry R.
Dwek. 2017. {``Imaging of {Pediatric Growth Plate Disturbances}.''}
\emph{RadioGraphics} 37 (6): 1791--1812.
\url{https://doi.org/10.1148/rg.2017170029}.

\bibitem[\citeproctext]{ref-ogdenDiagnosticImaging2000}
Ogden, John A., R. N. Hensinger, and N. McCollough. 2000. {``Diagnostic
{Imaging}.''} In \emph{Skeletal Injury in the Child}, 3rd ed., 128.
Springer.

\bibitem[\citeproctext]{ref-peinadocortesMechanobiologicalModelEpiphysis2011}
Peinado Cortés, L. M., J. C. Vanegas Acosta, and D. A. Garzón Alvarado.
2011. {``A Mechanobiological Model of Epiphysis Structures Formation.''}
\emph{Journal of Theoretical Biology} 287: 13--25.
\url{https://doi.org/10.1016/j.jtbi.2011.07.011}.

\bibitem[\citeproctext]{ref-sadeghianComputationalModelEndochondral2021}
Sadeghian, S. Mahsa, Frederic D. Shapiro, and Sandra J. Shefelbine.
2021. {``Computational Model of Endochondral Ossification: {Simulating}
Growth of a Long Bone.''} \emph{Bone} 153: 116132.
\url{https://doi.org/10.1016/j.bone.2021.116132}.

\bibitem[\citeproctext]{ref-stevensComputerModelEndochondral1999}
Stevens, S. S., G. S. Beaupré, and D. R. Carter. 1999. {``Computer Model
of Endochondral Growth and Ossification in Long Bones: Biological and
Mechanobiological Influences.''} \emph{Journal of Orthopaedic Research :
Official Publication of the Orthopaedic Research Society} 17 (5):
646--53. \url{https://doi.org/10.1002/jor.1100170505}.

\bibitem[\citeproctext]{ref-stoopTimingRestingZone2025}
Stoop, Jorik, Yuka Yokoyama, and Taiji Adachi. 2025. {``Timing of
Resting Zone Parathyroid Hormone-Related Protein Expression Affects
Maintenance of the Growth Plate During Secondary Ossification: A
Computational Study.''} \emph{Biomechanics and Modeling in
Mechanobiology} 24 (1): 125--37.
\url{https://doi.org/10.1007/s10237-024-01899-3}.

\bibitem[\citeproctext]{ref-sundaramurthyModulationEndochondralDevelopment2006}
Sundaramurthy, Sona, and Jeremy J. Mao. 2006. {``Modulation of
Endochondral Development of the Distal Femoral Condyle by Mechanical
Loading.''} \emph{Journal of Orthopaedic Research : Official Publication
of the Orthopaedic Research Society} 24 (2): 229--41.
\url{https://doi.org/10.1002/jor.20024}.

\bibitem[\citeproctext]{ref-vaca-gonzalezMechanobiologicalModelingEndochondral2018}
Vaca-González, Juan Jairo, Miguelangel Moncayo-Donoso, Johana Guevara,
Yoshie Hata, Sandra J. Shefelbine, and Diego Alexander Garzón-Alvarado.
2018. {``Mechanobiological Modeling of Endochondral Ossification: An
Experimental and Computational Analysis.''} \emph{Biomechanics and
Modeling in Mechanobiology} 17 (3): 853--75.
\url{https://doi.org/10.1007/s10237-017-0997-0}.

\bibitem[\citeproctext]{ref-wongTheoreticalModelEndochondral1990}
Wong, M., and D. R. Carter. 1990. {``A Theoretical Model of Endochondral
Ossification and Bone Architectural Construction in Long Bone
Ontogeny.''} \emph{Anatomy and Embryology} 181 (6): 523--32.
\url{https://doi.org/10.1007/BF00174625}.

\end{CSLReferences}




\end{document}
